% concatenated from main.tex
%%%%%%%%%%%%%%%%%%%%%%%%%%%%%%%%%%%%%%%%%%%%%%%%%%%%%%%%%%%%%%%%%%%%%%%%%%%%%%%%%%%%%%%%%%%%%%%%%%%
%%      This is a LaTeX file of the paper
%%
%%      Title: TBA
%%
%%      Maria Andrade
%%      Department of Mathematics
%%      Federal University of Sergipe
%%      49107-230, Sao Cristov\~ao-SE, Brazil
%%      maria@mat.ufs.br

%%     Allan Freitas
%%     Universidade Federal da Paraíba
%%    Departamento de Matemática
%%    João Pessoa - PB, Brazil
%%    allan@mat.ufpb.br

\documentclass[reqno,11pt]{amsart} 

\usepackage[colorlinks=true,urlcolor=blue,
citecolor=red,linkcolor=blue,linktocpage,pdfpagelabels,
bookmarksnumbered,bookmarksopen]{hyperref}
% \usepackage[notref,notcite]{showkeys}

\usepackage[left=2.6cm,right=2.6cm,top=2.9cm,bottom=2.9cm]{geometry}

\usepackage{amsmath,amssymb,latexsym,soul,cite,mathrsfs}
\numberwithin{equation}{section}  % Equation within section - Works with amsmath
\usepackage{color,enumitem,graphicx}
\usepackage{enumerate}
\usepackage{cancel}
%\usepackage{showlabels}
\usepackage{mathrsfs}
%\usepackage[notref,notcite]{showkeys}

\usepackage[colorlinks=true,urlcolor=blue,
citecolor=red,linkcolor=blue,linktocpage,pdfpagelabels,
bookmarksnumbered,bookmarksopen]{hyperref}
\usepackage[english]{babel}

\newtheorem{theorem}{Theorem}[section]
\newtheorem{proposition}[theorem]{Proposition}
\newtheorem{corollary}[theorem]{Corollary}
\newtheorem{lemma}[theorem]{Lemma}
\DeclareMathAlphabet{\mathpzc}{OT1}{pzc}{m}{it}
\newtheorem{definition}[theorem]{Definition}
\newtheorem{example}[theorem]{Example}
\newtheorem{xca}[theorem]{Exercise}
\newtheorem{remark}[theorem]{Remark}
\newcommand{\Ric}{\operatorname{Ric}}
\newtheorem{theoremletter}{Theorem}
\renewcommand{\thetheoremletter}{\Alph{theoremletter}}
\newcommand{\tr}{\text{tr}}
\renewcommand{\div}{\text{div}}
\newcommand{\tcr}{\textcolor{red}}
\newenvironment{tcb}
{\color{blue}}
{}%\textheight=25cm
 %\hoffset=-60pt
 %\voffset=-70pt
 %\parskip=1pt

%opening
\title[Torsional Rigidity and Spherical Deficit for a Dirichlet Problem]{Torsional Rigidity and Spherical Deficit for a Dirichlet Problem on Riemannian Manifolds}

\author[M. Andrade]{Maria Andrade}
\address[M. Andrade]{Universidade Federal de Sergipe - UFS, Departamento de Matemática, 49100-000, S\~ao Cristov\~ao-SE, Brazil.
}
\email{\href{mailto: maria@mat.ufs.br}{ maria@mat.ufs.br}}

\address{Current address: Department of Mathematics, Princeton University, Princeton, NJ, USA, 08544.}
\email{\href{mailto: ma6208@princeton.edu}{ ma6208@princeton.edu}}

\author[A. Freitas]{Allan Freitas}

\address[A. Freitas]{Universidade Federal da Paraíba - UFPB, Departamento de Matemática, João Pessoa, PB 58051-900, Brazil}
\email{\href{mailto: allan@mat.ufpb.br}{allan@mat.ufpb.br}}

\subjclass[2020]{53C20, 53C21, 58J32, 53C24}
%58J32 - Boundary value problems on manifolds
%53C18   	Conformal structures on manifolds
%53C20   	Global Riemannian geometry, including pinching
 %53C21   	Methods of global Riemannian geometry, including PDE methods; curvature restrictions 
\keywords{Torsional Rigidity, Ricci curvature, Overdetermined PDE, Model Manifold, Rigidity}


\begin{document}

\begin{abstract}
In this work, we study several inequalities related to a Dirichlet problem on Riemannian manifolds whose Ricci curvature is bounded from below. First, we establish inequalities involving the torsional rigidity and discuss rigidity results characterizing metric balls in this setting. Next, we derive an integral identity associated with a Dirichlet problem, which measures the spherical deficit arising in this context. In particular, we apply this identity to the setting of Einstein manifolds.
\end{abstract}

\maketitle

% Tarefas:

% Allan --- Termina introdução, mexer no início da Seção 4, escreve o último Corolário. Citar o Teorema 3.3 no caso de espaços modelo. Fazer abstract.
% \textcolor{red}{Allan, por favor, veja esta observação que você escreveu na seção 4 Pensar aqui se vale a pena colocar uma subseção abordando todos os resultados no setting acima. Pode ficar repetitivo fazer de um por um. Os resultados nesse setting podem ser interessantes para motivar a introdução.}

% Maria ---- Revisar todas as preliminares (colocar resultados de campos conformes), todas contas da Seção 4 e reesecrever a prova do Teorema 3.3. Conferir notações: $Vol(\Omega)$ ou $V(\Omega)$? Decidimos $u_{\nu\nu}$. Decidimos colocar $P$ (substituir $P(u)$) \tcr{feito, porém será revisto novamente.}




\section{Introduction}\label{sec:intro}
Let $\Omega$ be an bounded domain in the Euclidean space $\mathbb{R}^n$ with non-empty boundary $\partial\Omega$, and let the \textit{torsion function} $u : \Omega \to \mathbb{R}$ be the unique solution of
\begin{equation}\label{Dirichlet1}
\Delta u = -1\quad \text{in } \Omega, \quad u = 0 \text{ on } \partial \Omega.    
\end{equation}
In this setting, the \textit{torsional rigidity} is the functional defined by
\begin{equation}
T(\Omega) = \int_{\Omega} u \, dv.  
\end{equation}
The importance, applications, and meaning of both the torsion function and the torsional rigidity can be seen in many areas of mathematics, such as elasticity theory (see \cite{polya1951,sokolniko1956}), fluid dynamics (see \cite{serrin1971}), heat conduction (see \cite{vandenberg1990}), and also in the study of minimal submanifolds \cite{markvorsen2006}.
Just to mention an example, in the $n=2$ case, $T(\Omega)$ represents the torque required for a unit angle of twist per unit
length when twisting an elastic beam of uniform cross section $\Omega$ (see \cite[Pages~109-119]{sokolniko1956}).

In the case where an additional boundary condition is imposed together with \eqref{Dirichlet1}, namely
\begin{equation}\label{neumann}
u_{\nu} = -c,
\end{equation}
where $c$ is a positive constant, we are faced with a so-called overdetermined problem. This type of problem was first studied by J. Serrin in \cite{serrin1971}, who proved that the only solution of \eqref{Dirichlet1}--\eqref{neumann} is an Euclidean ball, and consequently the torsion function $u$ is radial. In an interesting application to elasticity theory, Serrin's celebrated result can be paraphrased as follows \cite{serrin1971}: \textit{``when a solid straight bar is subject to torsion, the magnitude of the resulting traction which occurs at the surface of the bar is independent of position if and only if the bar has a circular cross section.''}

In both the Euclidean setting, which has been primarily studied, and in the context of domains in \textit{Riemannian manifolds}, the analysis of problems of the type \eqref{Dirichlet1} has attracted considerable attention. A particularly interesting and recently explored direction is to establish bounds for the torsional rigidity $T(\Omega)$. In this regard, Wang and Xia (see \cite{wang2020}) have recently obtained sharp estimates for the torsional rigidity where $\Omega=(M^n,g)$ is a compact Riemannian manifold with boundary. In one of their results they proved 
\begin{theorem}[Theorem 1.7 in \cite{wang2020}]
Let $(M^n,g)$ be an $n$-dimensional compact Riemannian manifold with boundary and Ricci curvature $Ric\geq(n-1)k, \ k\in\mathbb{R}$. Let $u$ be the solution of the Dirichlet problem \eqref{Dirichlet1}.
Then
\begin{equation}\label{bound}
\max_{\partial M} |\nabla u|^2 
\geq 
\frac{(n+2)T(M)}{nVol(M)}
+ \frac{2(n-1)k}{Vol(M)} \int_M u\,|\nabla u|^2dv,
\end{equation}
with equality holding if and only if $\kappa = 0$ and $M$ is isometric to a ball in $\mathbb{R}^n$.
\end{theorem}    
In the Euclidean case, and in the presence of the boundary condition \eqref{neumann}, the previous result provides an alternative proof of Serrin's classical theorem. Indeed, in this particular case, equality in \eqref{bound} follows from the classical Pohozaev identity \cite{pohozaev1965} (see also \cite[p.~171]{struwe}), and the corresponding rigidity result then follows (see \cite{wang2020} for details). 

In the context of Riemannian manifolds, particularly those whose Ricci curvature satisfies $\mathrm{Ric} \geq (n-1)k g, \ k\in\mathbb{R},$ the following Dirichlet problem has attracted considerable attention in recent years:
\begin{equation}\label{Dirichlet2}
\Delta u + n k u = -1 \quad \text{in } \Omega, 
\qquad 
u = 0 \quad \text{on } \partial \Omega.
\end{equation}
This problem is typically studied in conjunction with the overdetermined condition \eqref{neumann}. In the case where $k$ corresponds to the sectional curvature of an $n$-dimensional space form, Ciraolo and Vezzoni \cite{ciraolo} extended Weinberger's method and established the corresponding rigidity result in $\mathbb{S}^{n}_{+}$ (when $k=1$) and $\mathbb{H}^{n}$ (when $k=-1$) by means of a suitable $P$-function. Subsequently, under the same lower bound on the Ricci curvature, this rigidity phenomenon has also been observed in certain classes of warped products (see \cite{Roncoroni2018, farina2022}) and in manifolds endowed with a conformal vector field (see \cite{freitas2024, andrade2025}).

The first goal of this paper is to establish new estimates for the torsional rigidity associated with problem \eqref{Dirichlet2}. To this end, we employ the classical $P$-function in this setting (see Section~\ref{sec_2}), extending the approach developed in \cite{wang2020} and obtaining new rigidity characterizations of metric balls in Riemannian manifolds. Although we obtain some results applicable to Riemannian manifolds with Ricci curvature bounded below (see Theorems \ref{thm:main} and \ref{thm:gradient_estimate}), we highlight one application in the context of so-called model manifolds.

Recall that a Riemannian manifold $(M_{\rho}^{n},g_{M_{\rho}^{n}})$ is called a \emph{model manifold} (or rotationally symmetric) if
\[
M_{\rho}^{n} \doteq \frac{[0,r)\times\mathbb{S}^{n-1}}{\sim}\qquad\text{and}\qquad 
g_{M_{\rho}^{n}} = dt\otimes dt + \rho^{2}(t)\,g_{\mathbb{S}^{n-1}},
\]
where $r\in(0,+\infty]$, $\sim$ is the equivalence relation identifying all points in $\{0\}\times\mathbb{S}^{n-1}$, and $\rho:[0,r)\to[0,+\infty)$ is a smooth function satisfying
\[
\rho(t)>0\ \ \forall t>0,\qquad \rho^{(2k)}(0)=0\ \ \forall k=0,1,2,\dots,\qquad \rho'(0)=1.
\]
The unique point corresponding to $t=0$ is called the \emph{pole} of the model and is denoted by $o\in M$, while $\rho$ is called the \emph{warping function} (see \cite[Section 1.4.4]{petersen}). Classical examples include the space forms: the Euclidean space $\mathbb{R}^n$, obtained by taking $\rho(t)=t$ and $r=\infty$; the hyperbolic space $\mathbb{H}^n$, obtained by taking $\rho(t)=\sinh(t)$ and $r=\infty$; and the upper hemisphere $\mathbb{S}_{+}^n$, obtained by taking $\rho(t)=\sin(t)$ and $r=\frac{\pi}{2}$. These three models serve as the prototypes for the following result, which is parallel to \cite[Theorem 1.7]{wang2020}:

\begin{theoremletter}\label{MainTh_1}
Let $(M_{\rho}^{n},g_{M_{\rho}^{n}})$ be a model manifold whose Ricci curvature satisfies $\operatorname{Ric} \geq (n-1)k g, \ k\in\mathbb{R}$ and let $\Omega \subset M_{\rho}^{n}$ be a bounded domain for which there exists a solution $u$ of \eqref{Dirichlet2}. Then
\begin{equation}\label{eqtor}
\max_{\partial\Omega}|\nabla u|^2 \;\geq\; \frac{1}{\operatorname{Vol}(\Omega) + nk\,T(\Omega)}\;\cdot\;\frac{n+2}{n}\left[ T(\Omega) + kn\int_{\Omega}\bigl(u^2 + u|\nabla u|^2\bigr)dv \right].
\end{equation}
Moreover, equality holds in \eqref{eqtor} if and only if $u$ is a radial function and $\Omega$ is a geodesic ball in a space form.
\end{theoremletter}

The second part of this paper is devoted to deriving an integral identity in order to analyze the spherical deficit associated with the solutions of \eqref{Dirichlet2}. The motivation for this goal comes from \cite{poggesi,poggesi1}, where the authors established new integral identities yielding quantitative estimates for radially symmetric configurations associated with the problem \eqref{Dirichlet1}. In this direction, for both Alexandrov's Soap Bubble Theorem and Serrin's overdetermined problem, the quantity
\[
|\mathring{\nabla}^{2}u|^{2}
=
|\nabla^{2}u|^{2}
-
\frac{(\Delta u)^{2}}{n}
\]
plays a fundamental role in measuring how close a domain is to being spherical, and for this reason we refer to it as the \textit{spherical (or Cauchy--Schwarz) deficit}. Indeed, the condition
\[
\mathring{\nabla}^{2}u=0
\]
leads to rigidity phenomena both in the Euclidean and in the Riemannian settings. In the Euclidean case, it characterizes quadratic polynomials, while in the Riemannian setting, Obata-type theorems identify metric balls as the only possible solutions (see, for example, \cite{obata, reilly} and \cite[Lemma 6]{farina2022}). In the next result, $q$ is a suitable quadratic polynomial solution for \eqref{Dirichlet1} and
$$c=\frac{\operatorname{Vol}(\Omega)}{\operatorname{Area}(\Omega)}.$$

\begin{theorem}[Theorem 2.1 in \cite{poggesi}]
Let $\Omega \subset \mathbb{R}^{n}$ be a bounded domain with boundary $\Gamma$ of class $C^{1,\alpha}$, with $0<\alpha\leq 1$, and let $c$ a positive constant. Then the solution $u$ of \eqref{Dirichlet1} satisfies the identity
\begin{equation}\label{eq:serrin_identity}
-2\int_{\Omega} u|\mathring{\nabla}^{2}u|^2dv
=\int_{\partial\Omega}(u_{\nu}^{2}-c^{2})(u_{\nu}-q_{\nu})d\sigma.
\end{equation}
In particular, if the right-hand side of \eqref{eq:serrin_identity} is non-positive, then $\partial\Omega$ must be a sphere (and hence $\Omega$ a ball) of radius $R=nc$. The same conclusion clearly holds whenever $u_{\nu}$ is constant on~$\partial\Omega$.
\end{theorem}

The main feature of this type of identity is that a boundary condition, namely, the constancy of the normal derivative of $u$, yields rigidity of the domain $\Omega$. Our second main result is devoted to establishing an integral identity in this direction for general Riemannian manifolds endowed with a conformal vector field (see Section~\ref{sec_2} for further details). Although we are able to derive a general integral identity for solutions of \eqref{Dirichlet2}, as stated in Proposition~\ref{intlemma}, which allows us to analyze the role played by the Ricci and scalar curvatures in these identities, we emphasize here its particularly intriguing application to the setting of Einstein manifolds, which we describe below. In this special case, the existence of solutions is guaranteed, for example, by \cite[Theorem 2.3]{foga} (see Remark \ref{existence} for details). In what follows, \(\varphi\) denotes a fundamental solution of the equation
\[
\mathring{\nabla}^{2}\varphi=0.
\]

\begin{theoremletter}\label{MainTh_2}
Let $(M^n,g)$ be an Einstein Riemannian manifold with $\Ric=(n-1)kg, \ k\in\mathbb{R}$, endowed with a conformal vector field $X$ with conformal factor $V$, and let $\varphi$ be a fundamental solution. If $u$ is a solution of \eqref{Dirichlet2} and $c$ is an arbitrary constant, then
\begin{align*}
-2\int_{\Omega}Vu\,|\mathring{\nabla}^{2}u|^{2}dv
&=
\int_{\partial\Omega}(u_{\nu}^{2}-c^{2})
\bigl[V(u-\varphi)_{\nu}-(u-\varphi)V_{\nu}\bigr]d\sigma. \\
% &\quad
% -2\int_{\Omega}u(u-\varphi)\langle \nabla^{2}u,\mathring{\nabla}^{2}V\rangle,
\end{align*}
% where $c$ is an arbitrary constant. In particular,
% \begin{itemize}
%     \item[(i)] if $u\leq \varphi$, then
%     \[
%     -2\int_{\Omega}Vu\,|\mathring{\nabla}^{2}u|^{2}
%     \leq
%     \int_{\partial\Omega}(u_{\nu}^{2}-c^{2})
%     \bigl[V(u-\varphi)_{\nu}-(u-\varphi)V_{\nu}\bigr];
%     \]
    
%     \item[(ii)] if $u\geq \varphi$, then
%     \[
%     -2\int_{\Omega}Vu\,|\mathring{\nabla}^{2}u|^{2}
%     \geq
%     \int_{\partial\Omega}(u_{\nu}^{2}-c^{2})
%     \bigl[V(u-\varphi)_{\nu}-(u-\varphi)V_{\nu}\bigr].
%     \]
% \end{itemize}
\end{theoremletter}

\begin{remark}
In the Euclidean setting, where the position vector field $x$ is conformal with conformal factor $V\equiv 1$, the previous result reduces to the classical identity established in \cite[Theorem 2.1]{poggesi}.
\end{remark}

\textbf{Organization of the paper.}
In Section~\ref{sec_2}, we recall some basic concepts and preliminary results that will be used throughout the paper. In Section~\ref{sec_estimates}, we investigate results related to torsional rigidity, establishing bounds and rigidity statements. In particular, we prove Theorem~\ref{MainTh_1}. Finally, in Section~\ref{sec_deficit}, we establish a general integral identity associated with problem \eqref{Dirichlet2}, extending the identity obtained in \cite[Theorem 2.1]{poggesi}; see Proposition~\ref{intlemma}. As an application, we prove Theorem~\ref{MainTh_2}.



\section{Preliminaries}\label{sec_2}

In this section, we gather the essential concepts and auxiliary results that will be used throughout the paper. Whenever appropriate, we highlight how each tool connects to the problems and goals outlined in the Introduction.

We begin by recalling Serrin's type problem in the Riemannian setting:
\begin{eqnarray}
\Delta u + nku  &=& -1 \quad \text{in} \quad \Omega, \quad u = 0 \quad \text{on} \quad \partial \Omega, \label{serrinoka}\\
u_{\nu} &=& -c, \quad \text{on} \quad \partial \Omega, \label{serrinokb}    
\end{eqnarray}
where $c$ is a constant and $k\in\mathbb{R}$. As discussed in the Introduction, the parameter $k$ reflects a lower bound on the Ricci curvature of the ambient manifold, and the overdetermined condition \eqref{serrinokb} is the distinctive feature of Serrin-type rigidity phenomena.

A key quantity associated with problem \eqref{serrinoka} is the torsional rigidity of $\Omega$. By integration by parts, one easily verifies that
\[
\int_{\Omega} |\nabla u|^2dv = \int_{\Omega} (-\Delta u)  u \, dv = \int_{\Omega} (1 + n k u)  u \, dv= \int_{\Omega}  u \, dv + n k \int_{\Omega} u^2 \,dv.
\]
This motivates the definition of torsional rigidity
\begin{equation}\label{eqtorsion}
T(\Omega) := \int_{\Omega} |\nabla u|^2\, dv - n k \int_{\Omega} u^2\, dv = \int_{\Omega} u \, dv.
\end{equation}
The quantity $T(\Omega)$ will be central in Section~\ref{sec_estimates}, where we establish sharp bounds for it under curvature assumptions.

A classical approach to Serrin's problem, both in Euclidean space and on Riemannian manifolds, employs an auxiliary function known as the $P$-function. In our setting, it is given by
\begin{equation}\label{pfunction}
P := |\nabla u|^2 + \dfrac{2}{n}u + ku^2.
\end{equation}
The importance of $P$ stems from the following fact: under the curvature assumption $\operatorname{Ric} \geq (n-1)k g$, the function $P$ is subharmonic for any solution of \eqref{serrinoka}. This was proved by Farina and Roncoroni \cite{farina2022} using the Bochner identity (see also \cite{ciraolo}), and its link with a Pozozaev-type identity is the cornerstone of the rigidity phenomenum in this setting.

More precisely, they established the following Obata-type result, which will be instrumental in the proofs of Theorems~\ref{thm:main} and \ref{thm:gradient_estimate}.

\begin{lemma}[\cite{farina2022}]\label{farina-roncoroni}
Let $(M^{n},g)$ be an $n$-dimensional Riemannian manifold satisfying  
\begin{equation}\label{ricci_bound}
\operatorname{Ric} \geq (n-1)k g, \quad k \in \mathbb{R}.    
\end{equation}
Let $\Omega \subset M$ be a domain and $u \in C^{2}(\Omega)$ a solution of 
\[
\Delta u + nku = -1 \quad \text{in } \Omega.
\]
Then the $P$-function defined in \eqref{pfunction} satisfies 
\[
\Delta P \geq 0.
\]
Moreover, $\Delta P= 0$ if and only if $\Omega$ is a metric ball and $u$ is a radial function, i.e., $u$ depends only on the distance from the center of the ball. In this case, $u$ also satisfies
\[
\nabla^2 u = -\left(\frac{1}{n} + ku\right)g \quad \text{in } \Omega,
\]
and $\operatorname{Ric}(\nabla u, \nabla u) = (n-1)k|\nabla u|^2$ in $\Omega$.
\end{lemma}

In Section~\ref{sec_deficit}, we shift our perspective to study the spherical deficit for solutions of \eqref{serrinoka}. This requires a richer geometric setting: Riemannian manifolds endowed with a conformal vector field. Recall that a vector field $X$ on $(M^n,g)$ is called conformal if its flow consists of conformal transformations. Equivalently, there exists a smooth function $V$ on $M$ such that the Lie derivative of the metric satisfies $\mathcal{L}_X g = 2V g$. The function $V$ is called the conformal factor of $X$. When $X$ is also closed (i.e., its dual $1$-form is closed), we obtain a particularly tractable structure that generalizes the position vector field in Euclidean space.

The following lemma, which is treated for example in \cite[Lemma 2.2]{andrade2025}, relates the conformal factor $V$ to the scalar curvature $R$ of the manifold. It will be crucial in the derivation of our general integral identity (Proposition~\ref{intlemma}).

\begin{lemma}\label{lemma1}
Let $(M^{n},g)$ be a Riemannian manifold endowed with a conformal vector field $X$ with conformal factor $V$. Then
\begin{equation}\label{Lapl_V}
-(n-1)\Delta V
=
\frac{1}{2}X(R)+VR,
\end{equation}
where $R$ denotes the scalar curvature of $(M^{n},g)$.
\end{lemma}


We also recall the following result, particularly applicable in the case of Einstein manifolds. It is a consequence of \cite[Lemma 2.1]{kr}, see also \cite[p.~160]{yano}.


\begin{lemma}\label{lemma_tracefree}
Let $(M^{n},g)$ be a Riemannian manifold endowed with a conformal vector field $X$ with conformal factor $V$. Then
\[
\mathcal{L}_{X}\mathring{\operatorname{Ric}}=0
\quad\Longleftrightarrow\quad
\mathring{\nabla}^{2}V=0.
\]

In particular, if $(M^{n},g)$ is Einstein, then
\[
\mathring{\nabla}^{2}V=0.
\]
\end{lemma}

\begin{remark}
We recall that a conformal vector field $X$ is said to be \emph{closed conformal} if the $1$-form dual to $X$ is closed. Equivalently,
\[
\nabla_YX = VY,
\]
for every $Y\in\mathfrak{X}(M)$, where $V$ denotes the conformal factor of $X$. A remarkable property of manifolds admitting a nontrivial closed conformal vector field is that they are locally isometric to a warped product with a one-dimensional factor; see, for instance, \cite[Section 3]{montiel}. This class includes space forms and many other geometrically relevant examples.
\end{remark}

% As motivated in the Introduction, the quantity
% \[
% |\mathring{\nabla}^{2}u|^{2} := |\nabla^{2}u|^{2} - \frac{(\Delta u)^{2}}{n}
% \]
% measures the deviation of the Hessian of $u$ from being purely trace-like, and is therefore called the spherical (or Cauchy--Schwarz) deficit. In the Euclidean case, the condition $\mathring{\nabla}^{2}u = 0$ characterizes quadratic polynomials, while in the Riemannian setting it leads, via Obata-type theorems, to the rigidity of metric balls.

% For our analysis in Section~\ref{sec_deficit}, we also need a fundamental solution $\varphi$ of the linearized problem, defined by the condition $\mathring{\nabla}^{2}\varphi = 0$, which in the context of problem \eqref{serrinoka} is equivalent to requiring
% \[
% \nabla^2 \varphi = -\left(\frac{1}{n} + k\varphi\right)g,
% \]
% and consequently $\Delta \varphi = -(1 + nk\varphi)$. In the Euclidean case ($k=0$), such a $\varphi$ is a quadratic polynomial, and the difference $u-\varphi$ appears naturally in integral identities that control the spherical deficit (see \cite{poggesi} and Lemma~\ref{intlemma} below). This idea will be extended to the Riemannian setting in Section 4, where we derive a general integral identity (Lemma \ref{intlemma}) that controls the spherical deficit through the difference $u - \varphi$.

\section{Estimates of the Torsional Rigidity}\label{sec_estimates}

We start this section giving our first result. Such a result may be compared with \cite[Theorem 1.6]{wang2020}.

\begin{theorem}\label{thm:main}
Let $(M^n,g)$ be an $n$-dimensional Riemannian manifold whose Ricci curvature satisfies $\operatorname{Ric}\geq (n-1)kg$, $k\in\mathbb{R}$, and let $\Omega\subset M$ be a bounded domain for which there exists a solution $u$ of
\begin{equation}\label{eq:main_problem}
\Delta u + nku = -1\quad \text{in } \Omega,\qquad u = 0\quad \text{on } \partial\Omega.
\end{equation}
Let $\nu$ denote the outward unit normal vector field along $\partial\Omega$.

\begin{itemize}
\item[(a)] We have
\begin{equation}\label{eq:min_unn}
\min_{x\in\partial\Omega} u_{\nu\nu}(x)\leq -\frac{1}{n}.
\end{equation}
Moreover, equality holds in \eqref{eq:min_unn} if and only if $u$ is radial and $\Omega$ is a metric ball in $(M^n,g)$.

\item[(b)] Let $\operatorname{Area}(\partial\Omega)$, $H$, $T(\Omega)$, and $\operatorname{Vol}(\Omega)$ denote the area of $\partial\Omega$, the mean curvature of $\partial\Omega$, the torsional rigidity of $\Omega$, and the volume of $\Omega$, respectively. If $H\geq 0$ on $\partial\Omega$, then
\begin{equation}\label{eq:integral_unn}
\int_{\partial\Omega}u_{\nu\nu}d\sigma\leq \left(\frac{n-1}{n}\right)^{1/2}\left(\operatorname{Vol}(\Omega) + nkT(\Omega)\right)^{1/2}\left(\int_{\partial\Omega}Hd\sigma\right)^{1/2} - \operatorname{Area}(\partial\Omega).
\end{equation}
Equality holds in \eqref{eq:integral_unn} if and only if $u$ is radial and $\Omega$ is a metric ball in $(M^n,g)$.
\end{itemize}
\end{theorem}

\begin{proof}
Since $u$ satisfies \eqref{eq:main_problem}, we have $\Delta u<0$. Although this is immediate when $k\geq 0$, it also remains true in the case $k<0$ by \cite[Lemma 5.2]{andrade2025}. Therefore, the strong maximum principle together with Hopf's lemma yields
\begin{equation}\label{eq:unu_negative}
u_{\nu}(x)<0,\qquad \forall x\in\partial\Omega.
\end{equation}

Using the classical Bochner formula
\[\frac{1}{2}\Delta |\nabla u|^2 = |\nabla^2 u|^2 +\langle \nabla u,\nabla \Delta u\rangle +\operatorname{Ric}(\nabla u,\nabla u),\]
together with \eqref{eq:main_problem} and the assumption $\operatorname{Ric}\geq (n-1)kg$, we obtain
\begin{equation}\label{eq:bochner_estimate}
\frac{1}{2}\Delta |\nabla u|^2\geq |\nabla^2 u|^2 -nk|\nabla u|^2 +(n-1)k|\nabla u|^2 = |\nabla^2 u|^2 -k|\nabla u|^2.
\end{equation}
Integrating over $\Omega$ and applying the divergence theorem, we get
\begin{equation}\label{eq:integral_bochner}
\int_{\Omega}|\nabla^2 u|^2dv -k\int_{\Omega}|\nabla u|^2dv\leq \frac{1}{2}\int_{\partial\Omega}\frac{\partial}{\partial\nu} |\nabla u|^2d\sigma.
\end{equation}
Since $u = 0$ on $\partial\Omega$, it follows that
\begin{equation}\label{eq:boundary_grad}
\nabla u = u_{\nu}\nu\ \text{on }\partial\Omega,\quad\text{and}\quad \frac{\partial}{\partial\nu} |\nabla u|^2 = 2\nabla^2 u(\nabla u,\nu) = 2u_{\nu}u_{\nu\nu}\ \text{on }\partial\Omega.
\end{equation}
Combining this with \eqref{eq:integral_bochner}, we conclude that
\begin{equation}\label{eq:inequality_hessian}
\int_{\Omega}|\nabla^2 u|^2dv -k\int_{\Omega}|\nabla u|^2dv\leq \int_{\partial\Omega}u_{\nu}u_{\nu\nu}d\sigma.
\end{equation}

Define
\[L = \min_{x\in\partial\Omega} u_{\nu\nu}(x).\]
Then, by \eqref{eq:unu_negative} and the divergence theorem, we obtain
\[\int_{\partial\Omega}u_{\nu}u_{\nu\nu}d\sigma\leq L\int_{\partial\Omega}u_{\nu}d\sigma = L\int_{\Omega}\Delta u dv= -L\int_{\Omega}(nku+1)dv.\]
Substituting this into \eqref{eq:inequality_hessian}, we infer
\begin{equation}\label{eq:after_L}
\int_{\Omega}|\nabla^2 u|^2dv -k\int_{\Omega}|\nabla u|^2dv\leq -nkL\int_{\Omega}u \, dv- L\operatorname{Vol}(\Omega).
\end{equation}
Using the definition of torsional rigidity $T(\Omega)=\int_{\Omega}udv$ and the identity $\int_{\Omega}|\nabla u|^2dv = T(\Omega)+nk\int_{\Omega}u^2dv$, we rewrite \eqref{eq:after_L} as
\begin{equation}\label{eq:before_CS}
\int_{\Omega}|\nabla^2 u|^2dv -nk^2\int_{\Omega}u^2dv +kT(\Omega)(-1+nL)\leq -L\operatorname{Vol}(\Omega).
\end{equation}

On the other hand, the Cauchy-Schwarz inequality together with \eqref{eq:main_problem} yields
\[|\nabla^2 u|^2\geq \frac{(\Delta u)^2}{n} = \frac{(1+nku)^2}{n} = nk^2 u^2 +2ku +\frac{1}{n}.\]
Integrating over $\Omega$, we obtain
\begin{equation}\label{eq:CS_integral}
\int_{\Omega}|\nabla^2 u|^2dv\geq nk^2\int_{\Omega}u^2\, dv +2k\int_{\Omega}u\, dv +\frac{1}{n}\operatorname{Vol}(\Omega).
\end{equation}
Substituting \eqref{eq:CS_integral} into \eqref{eq:before_CS}, we arrive at
\[L\left(nkT(\Omega) + \operatorname{Vol}(\Omega)\right)\leq -\frac{1}{n}\left(nkT(\Omega) + \operatorname{Vol}(\Omega)\right).\]
We note that $nkT(\Omega) + \operatorname{Vol}(\Omega) > 0$: this is immediate when $k\geq 0$, while for $k<0$ it follows from \cite[Lemma 5.2]{andrade2025}. Therefore,
\begin{equation}\label{eq:L_bound}
L\leq -\frac{1}{n}.
\end{equation}
This proves part (a) of the theorem.

We now analyze the equality case. Equality holds in \eqref{eq:L_bound} if and only if equality holds in \eqref{eq:bochner_estimate} and \eqref{eq:CS_integral}, which implies
\begin{equation}\label{eq:rigidity_conditions}
\nabla^2 u = \frac{\Delta u}{n} g\qquad\text{and}\qquad \operatorname{Ric}(\nabla u,\nabla u) = (n-1)k|\nabla u|^2.
\end{equation}
In particular, the first condition implies that $u$ is radial and that $\Omega$ is a metric ball, by an Obata-type result (see, e.g., \cite[Lemma 6]{farina2022}).

We now prove part (b). Since $u$ vanishes on $\partial\Omega$, it follows from \eqref{eq:main_problem} that $\Delta u = -1$ on $\partial\Omega$. On the other hand, the decomposition of the Laplacian along the boundary gives
\[\Delta u = \Delta_{\partial\Omega}u + Hu_{\nu} + u_{\nu\nu}.\]
Because $u = 0$ on $\partial\Omega$, we have $\Delta_{\partial\Omega}u = 0$. Hence, integrating over $\partial\Omega$, we deduce
\begin{equation}\label{eq:boundary_decomposition}
\operatorname{Area}(\partial\Omega) + \int_{\partial\Omega}u_{\nu\nu}d\sigma = -\int_{\partial\Omega}Hu_{\nu}d\sigma.
\end{equation}

We now invoke Reilly's formula \cite{reilly}, which in the present setting reads
\[\int_{\Omega}\left((\Delta u)^2 - |\nabla^2 u|^2 - \operatorname{Ric}(\nabla u,\nabla u)\right) dv= \int_{\partial\Omega}H u_{\nu}^2d\sigma.\]
Substituting $\Delta u = -(1 + nku)$ into this identity, we obtain
\[\int_{\partial\Omega}H u_{\nu}^2d\sigma = \int_{\Omega}\left((1 + nku)^2 - |\nabla^2 u|^2 - \operatorname{Ric}(\nabla u,\nabla u)\right)dv.\]
Using the estimates $|\nabla^2 u|^2\geq \frac{(\Delta u)^2}{n}$ and $\operatorname{Ric}(\nabla u,\nabla u)\geq (n-1)k|\nabla u|^2$, we infer
\begin{equation}\label{eq:reilly_estimate}
\int_{\partial\Omega}H u_{\nu}^2d\sigma\leq \frac{n-1}{n}\int_{\Omega}\left(1 + 2nku + n^2k^2u^2\right)dv - (n-1)k\int_{\Omega}|\nabla u|^2dv.
\end{equation}
Moreover, equality holds in \eqref{eq:reilly_estimate} if and only if \eqref{eq:rigidity_conditions} is satisfied.

Integration by parts together with the boundary condition $u = 0$ on $\partial\Omega$ gives
\[\int_{\Omega}|\nabla u|^2dv = -\int_{\Omega}u\Delta udv = \int_{\Omega}udv + nk\int_{\Omega}u^2dv = T(\Omega) + nk\int_{\Omega}u^2dv.\]
Combining this with \eqref{eq:reilly_estimate}, we arrive at
\begin{equation}\label{eq:H_bound}
\int_{\partial\Omega}H u_{\nu}^2d\sigma\leq \frac{n-1}{n}\operatorname{Vol}(\Omega) + (n-1)kT(\Omega).
\end{equation}
Since $H\geq 0$ on $\partial\Omega$ and $u_{\nu}<0$ by Hopf's lemma, H\"older's inequality yields
\[-\int_{\partial\Omega}Hu_{\nu}d\sigma\leq \left(\int_{\partial\Omega}H u_{\nu}^2d\sigma\right)^{1/2}\left(\int_{\partial\Omega}Hd\sigma\right)^{1/2}.\]
Substituting \eqref{eq:H_bound} into this estimate and then into \eqref{eq:boundary_decomposition}, we conclude
\[\int_{\partial\Omega}u_{\nu\nu}\leq \left(\frac{n-1}{n}\right)^{1/2}\left(\operatorname{Vol}(\Omega) + nkT(\Omega)\right)^{1/2}\left(\int_{\partial\Omega}Hd\sigma\right)^{1/2} - \operatorname{Area}(\partial\Omega),\]
which is exactly \eqref{eq:integral_unn}. Equality holds if and only if \eqref{eq:rigidity_conditions} holds, and the rigidity conclusion follows from \cite[Lemma 6]{farina2022}. This completes the proof of Theorem \ref{thm:main}.
\end{proof}

We now analyze the rigidity condition \eqref{eq:rigidity_conditions} in a fundamental case: the so-called model manifolds, as motivated in the Introduction. The space forms are the substrate of the following consequence.

% Recall that a Riemannian manifold $(M_{\rho}^{n},g_{M_{\rho}^{n}})$ is called a \emph{model manifold} if
% \[M_{\rho}^{n}\doteq \frac{[0,R)\times\mathbb{S}^{n-1}}{\sim}\quad\text{and}\quad g_{M_{\rho}^{n}}= dt\otimes dt + \rho^{2}(t)g_{\mathbb{S}^{n-1}},\]
% where $R\in (0,+\infty]$, $\sim$ is the equivalence relation that identifies all points in $\{0\}\times\mathbb{S}^{n-1}$, and $\rho:[0,R)\to[0,+\infty)$ is a smooth function satisfying:
% \[\rho(t)>0\ \forall t>0,\qquad \rho^{(2k)}(0)=0\ \forall k=0,1,2,\ldots,\qquad \rho'(0)=1.\]
% Moreover, the unique point corresponding to $t=0$ is called the pole of the model and is denoted by $o\in M$, while $\rho$ is called the warping function (see \cite{petersen}). Classical examples include the space forms: the Euclidean space $\mathbb{R}^n$, obtained by taking $\rho(t)=t$ and $R=\infty$; the hyperbolic space $\mathbb{H}^n$, obtained by taking $\rho(t)=\sinh(t)$ and $R=\infty$; and the upper hemisphere $\mathbb{S}_{+}^n$, obtained by taking $\rho(t)=\sin(t)$ and $R=\frac{\pi}{2}$.

\begin{corollary}\label{cor:model}
Let $(M_{\rho}^{n},g_{M_{\rho}^{n}})$ be an $n$-dimensional model manifold whose Ricci curvature satisfies $\operatorname{Ric}\geq (n-1)kg$, $k\in\mathbb{R}$, and let $\Omega\subset M$ be a bounded domain for which there exists a solution $u$ to \eqref{eq:main_problem}. Let $\nu$ denote the outward unit normal vector field along $\partial\Omega$.

\begin{itemize}
\item[(a)] We have
\[\min_{x\in\partial\Omega}u_{\nu\nu}(x)\leq -\frac{1}{n}.\]
Moreover, equality holds if and only if $\Omega$ is a geodesic ball in a space form and $u_{\nu\nu}=-\frac{1}{n}$ on $\partial\Omega$.

\item[(b)] Let $\operatorname{Area}(\partial\Omega)$, $H$, $T(\Omega)$, and $\operatorname{Vol}(\Omega)$ denote the area of $\partial\Omega$, the mean curvature of $\partial\Omega$, the torsional rigidity of $\Omega$, and the volume of $\Omega$, respectively. If $H\geq 0$ on $\partial\Omega$, then
\[\int_{\partial\Omega}u_{\nu\nu}\leq (n(n-1))^{1/2}\left(\operatorname{Vol}(\Omega) + nkT(\Omega)\right)^{1/2}\left(\int_{\partial\Omega}Hd\sigma\right)^{1/2} - \operatorname{Area}(\partial\Omega).\]
Equality holds if and only if $u$ is radial and $\Omega$ is a geodesic ball in a space form.
\end{itemize}
\end{corollary}

\begin{proof}
Besides the fact that $\Omega$ is a metric ball (which follows from the first condition in \eqref{eq:rigidity_conditions}), we now analyze more carefully the rigidity condition
\[\operatorname{Ric}(\nabla u,\nabla u) = (n-1)k|\nabla u|^2.\]
In the particular setting of model manifolds, since $(M,g)$ is rotationally symmetric and $u$ is radial, we have
\[\operatorname{Ric}(\nabla u,\nabla u) = -(n-1)\frac{\rho''(t)}{\rho(t)}|\nabla u|^2.\]
Consequently,
\[-\frac{\rho''(t)}{\rho(t)} = k,\qquad t\in(0,R).\]
Combining this equation with the smoothness conditions at the pole,
\[\rho^{(2k)}(0)=0,\qquad \rho'(0)=1,\]
we conclude that $\rho$ satisfies the differential equation $\rho'' + k\rho = 0$, whose solutions are
\[\rho(t) = 
\begin{cases}
\frac{1}{\sqrt{k}}\sin(\sqrt{k}\,t), & k>0,\\[4pt]
t, & k=0,\\[4pt]
\frac{1}{\sqrt{-k}}\sinh(\sqrt{-k}\,t), & k<0.
\end{cases}\]
These correspond exactly to the space forms $\mathbb{S}^n$, $\mathbb{R}^n$, and $\mathbb{H}^n$, respectively. Hence, the rigidity conditions force $\Omega$ to be a geodesic ball in a space form. The statements (a) and (b) then follow directly from Theorem \ref{thm:main}.
\end{proof}

Now, we explicitly give a limitation for the torsional rigidity in manifolds with Ricci bounded below. Such a result could be compared with \cite[Theorem 1.7]{wang2020}.

\begin{theorem}\label{thm:gradient_estimate}
Let $(M^n,g)$ be an $n$-dimensional Riemannian manifold whose Ricci curvature satisfies $\operatorname{Ric}\geq (n-1)kg$, $k\in\mathbb{R}$, and let $\Omega\subset M$ be a bounded domain for which there exists a solution $u$ of \eqref{eq:main_problem}. Then
\begin{equation}\label{eq:max_gradient_estimate}
\max_{\partial\Omega}|\nabla u|^2\geq \frac{1}{\operatorname{Vol}(\Omega) + nkT(\Omega)}\cdot\frac{n+2}{n}\left[T(\Omega) + nk\int_{\Omega}\left(u^2 + u|\nabla u|^2\right)dv\right].
\end{equation}
Moreover, equality holds in \eqref{eq:max_gradient_estimate} if and only if $u$ is a radial function and $\Omega$ is a metric ball in $(M^n,g)$.
\end{theorem}

\begin{proof}
Throughout the proof, we consider the auxiliary $P$-function
\[P = |\nabla u|^2 +\frac{2}{n}u + ku^2.\]
Using integration by parts, together with the boundary conditions $u = 0$ on $\partial\Omega$ and $P = |\nabla u|^2 = u_{\nu}^2$ on $\partial\Omega$, we obtain
\begin{eqnarray}\label{eq323}
\int_{\Omega}(u\Delta P - P\Delta u)dv = \int_{\partial\Omega}(uP_{\nu} - Pu_{\nu})d\sigma = -\int_{\partial\Omega}u_{\nu}^{3}d\sigma.
\end{eqnarray}
Recall that, under the assumption $\operatorname{Ric}\geq (n-1)kg$,
\[\Delta P = 2\Big(|\mathring{\nabla^2} u|^2 +(\operatorname{Ric} - (n-1)kg)(\nabla u,\nabla u)\Big)\geq 0.\]
Moreover, equality holds if and only if \eqref{eq:rigidity_conditions} is satisfied.

Let
\[m = \max_{\partial\Omega}|\nabla u|.\]
Substituting the above information into \eqref{eq323}, and using the fact that $u_{\nu}<0$ (see the beginning of the proof of Theorem \ref{thm:main}), we deduce
\[-\int_{\Omega}P\Delta udv\leq -m^2\int_{\partial\Omega}u_{\nu}d\sigma = m^2\int_{\Omega}(-\Delta u) dv= m^2\big(\operatorname{Vol}(\Omega) + nkT(\Omega)\big),\]
where we used $-\Delta u = 1 + nku$ and $\int_{\Omega}udv = T(\Omega)$.

Next, using the equation $-\Delta u = 1 + nku$, we compute
\begin{eqnarray*}
-\int_{\Omega}(\Delta u)P dv&=& \int_{\Omega}(1 + nku)\left(|\nabla u|^2 + \frac{2}{n}u + ku^2\right)dv\\
&=&\displaystyle\int_{\Omega}|\nabla u|^2dv + \frac{2}{n}\int_{\Omega}udv + k\int_{\Omega}u^2dv + nk\int_{\Omega}u|\nabla u|^2dv + 2k\int_{\Omega}u^2dv + nk^2\int_{\Omega}u^3dv.
\end{eqnarray*}

Using the identity $\int_{\Omega}|\nabla u|^2dv = \int_{\Omega}udv + nk\int_{\Omega}u^2dv$, it follows that
\[-\int_{\Omega}(\Delta u)P dv= \frac{n+2}{n}\int_{\Omega}udv + k(n+3)\int_{\Omega}u^2dv + nk\int_{\Omega}u|\nabla u|^2dv + nk^2\int_{\Omega}u^3dv.\]

On the other hand, integrating by parts and using $\operatorname{div}(u^2\nabla u) = u^2\Delta u + 2u|\nabla u|^2$, together with the boundary condition $u = 0$ on $\partial\Omega$, we obtain
\[nk^2\int_{\Omega}u^3dv = 2k\int_{\Omega}u|\nabla u|^2dv - k\int_{\Omega}u^2dv.\]

Therefore,
\[-\int_{\Omega}(\Delta u)Pdv = \frac{n+2}{n}\int_{\Omega}udv + k(n+2)\left(\int_{\Omega}u^2dv + \int_{\Omega}u|\nabla u|^2dv\right).\]

Finally, combining the estimates, we conclude that
\[m^2\big(\operatorname{Vol}(\Omega) + nkT(\Omega)\big)\geq \frac{n+2}{n}\left[T(\Omega) + nk\int_{\Omega}\left(u^2 + u|\nabla u|^2\right)dv\right].\]

Since $\operatorname{Vol}(\Omega) + nkT(\Omega) > 0$, which is immediate when $k\geq 0$, and follows from \cite[Lemma 5.2]{andrade2025} when $k<0$, the desired conclusion follows. Moreover, equality in the estimate implies that $\nabla^2 u = \frac{\Delta u}{n}g$, hence the rigidity conclusion follows from \cite[Lemma 6]{farina2022}.
\end{proof}

\begin{proof}[Proof of Theorem \ref{MainTh_1}]
It is sufficient to analyze the equality case in the Ricci curvature condition appearing in \eqref{eq:rigidity_conditions}. The conclusion then follows by arguing exactly as in Corollary \ref{cor:model}.
\end{proof}

\begin{remark}
 In the direction of \cite[Theorem 1.7]{wang2020}, if we start with a compact model Riemannian manifold with a boundary $(M_{\rho}^{n},g_{M_{\rho}^{n}})$, where there exists a solution for 
 \begin{equation}\label{Dirichlet3}
\Delta u + n k u = -1 \quad \text{in } M_{\rho}, 
\qquad 
u = 0 \quad \text{on } \partial M_{\rho},
\end{equation}
Corollary \ref{cor:model} and Theorem \ref{MainTh_1} implies that, the equality in those inequalities holds if and only if $(M_{\rho}^{n},g_{M_{\rho}^{n}})$ is isometric to a geodesic ball in a space form.
\end{remark}
%\section{tentativa} \textcolor{red}{sem ajustes, apenas contas}

%Pohozaev identity: $(M^n,g)$ Riemannian manifold, $X$ is a conformal vector field
%Está provada no pdf tablet

\section{An Integral identity and the spherical deficit}\label{sec_deficit}

Along this section, we assume that $(M^{n},g)$ is a Riemannian manifold endowed with a conformal vector field $X$. This means that the Lie derivative of the metric satisfies
\[
\mathcal{L}_{X}g = 2Vg,
\]
for some smooth function $V\colon M\to\mathbb{R}$, called the conformal factor of $X$. Although conformal vector fields naturally arise from one-parameter families of conformal transformations, it is worthwhile to mention some classical examples.

\begin{example}[Space forms]\label{space_forms}
The simply connected space forms admit natural conformal vector fields.

\begin{itemize}
    \item In the Euclidean space $\mathbb{R}^{n}$ endowed with the Euclidean metric $g_{0}$, the position vector field is conformal with conformal factor $V\equiv 1$.

    \item In the hyperbolic space $\mathbb{H}^{n}$ written as
    \[
    \mathbb{H}^{n}=([0,\infty)\times \mathbb{S}^{n-1},\,dr^{2}+\sinh^{2}(r)\,g_{\mathbb{S}^{n-1}}),
    \]
    the vector field \(X=\sinh(r)\,\partial_{r}\)
    is conformal with conformal factor $V=\cosh(r)$.
    

    \item In the sphere $\mathbb{S}^{n}$ written as
    \[
    \mathbb{S}^{n}=((0,\pi)\times \mathbb{S}^{n-1},\,dr^{2}+\sin^{2}(r)\,g_{\mathbb{S}^{n-1}}),
    \]
    the vector field $X=\sin(r)\,\partial_{r}$
    is conformal with conformal factor $V=\cos(r)$.
\end{itemize}
\end{example}

\begin{example}[Warped products]
More generally, let
\[
(M^{n},g)=\bigl((a,b)\times N^{n-1},\,dr^{2}+h(r)^{2}g_{N}\bigr)
\]
be a warped product manifold, where $(N^{n-1},g_{N})$ is a Riemannian manifold and $h\colon(a,b)\to\mathbb{R}_{+}$ is smooth. Then the vector field $X=h(r)\,\partial_{r}$ is conformal with conformal factor $V=h'(r)$.
\end{example}

Furthermore, we consider a fundamental solution $\varphi$ of \eqref{Dirichlet2}, namely a function satisfying
\begin{equation}\label{deltaphi1}
\nabla^{2}\varphi
=
-\left(\frac{1}{n}+K\varphi\right)g
\end{equation}
in $\Omega$. Taking the trace of \eqref{deltaphi1}, we obtain
\begin{equation}\label{deltaphi}
\Delta \varphi
=
-\left(1+nK\varphi\right).
\end{equation}

Before stating our first result, we present some examples of such solutions in particular settings.

\begin{example}[Euclidean space]
In the Euclidean case, where $K=0$, a corresponding fundamental solution is given by
\[
\varphi(x)=\frac{1}{2n}\bigl(a-|x-b|^{2}\bigr),
\]
where $a\in\mathbb{R}$ and $b\in\mathbb{R}^{n}$ are constants.
\end{example}

\begin{example}[Sphere and hyperbolic space]\label{ex:space_forms_phi}
Consider the space forms described in Example~\ref{space_forms}. Namely,
\begin{itemize}
    \item the hyperbolic space, for which
    \[
    h(r)=\sinh r
    \qquad \text{and} \qquad K=-1,
    \]
    
    \item and the sphere, for which
    \[
    h(r)=\sin r
    \qquad \text{and} \qquad K=1.
    \]
\end{itemize}

In both cases, the warping function satisfies
\[
h''=-Kh,
\]
and consequently the radial function
\[
\varphi(r)=c_{0}h'(r)-\frac{1}{nK},
\]
where $c_{0}$ is a constant, satisfies
\[
\nabla^{2}\varphi
=
-\left(\frac{1}{n}+K\varphi\right)g.
\]
\end{example}

% \begin{remark}
% We point out that manifolds endowed with a nontrivial closed conformal vector field possess an interesting property: they are locally isometric to a warped product with a 1-dimensional factor (for more details, see, e.g., \cite[Section 3]{montiel}).    
% \end{remark}


In the following, we prove an integral identity which constitutes the core of this section. As discussed in the introduction, it generalizes \cite[Theorem 2.1]{poggesi}, and the strength of identities of this type lies in their ability to measure the spherical deficit through the quantity
\[
\Delta P
=
2\left(
|\mathring{\nabla}^{2}u|^{2}
+
(\operatorname{Ric}-(n-1)kg)(\nabla u,\nabla u)
\right),
\]
which appears on the left-hand side (see Remark~\ref{remark_main} below).
\begin{proposition}[General Integral Identity]\label{intlemma}
Let $(M^n, g)$ be a Riemannian manifold endowed with a conformal vector field $X$ with conformal factor $V$. Let $\Omega$ be a bounded domain with $C^2$ boundary, and let $u$ be a solution of
\[
\Delta u=-(1+nku), \qquad k\in\mathbb{R},
\]
in $\Omega$. Assume further that $\varphi$ is a solution of \eqref{deltaphi1}. Then the following identity holds:
\begin{eqnarray}\label{eqlema}
\displaystyle\int_{\Omega}-Vu\Delta Pdv&=&\displaystyle\int_{\partial \Omega}|\nabla u|^2[V(u-\varphi)_{\nu}-(u-\varphi)V_{\nu}]d\sigma
+2\displaystyle\int_{\partial \Omega}u\nabla^2u((u-\varphi)\nabla V-V\nabla (u-\varphi),\nu)d\sigma\nonumber\\
&&+k\displaystyle\int_{\partial \Omega}u^2[(u-\varphi)V_{\nu}-V(u-\varphi)_{\nu}]d\sigma
-\dfrac{2}{n(n-1)}\displaystyle\int_{\Omega}Vu(u-\varphi)(R-n(n-1)k)dv\nonumber\\
&&-2\displaystyle\int_{\Omega}u(u-\varphi)(\operatorname{Ric}-(n-1)kg)(\nabla u,\nabla V)dv
-2\displaystyle\int_{\Omega}Vu(\operatorname{Ric}-(n-1)kg)(\nabla u,\nabla \varphi)dv\nonumber\\
&&-\dfrac{k}{n-1}\displaystyle\int_{\Omega}Vu^2(u-\varphi)(R-(n-1)nk)dv
-\dfrac{1}{n-1}\displaystyle\int_{\Omega}V|\nabla u|^2(u-\varphi)(R-(n-1)nk)dv\nonumber\\
&&-\dfrac{k}{2(n-1)}\displaystyle\int_{\Omega}u^2(u-\varphi)X(R)dv
-\dfrac{1}{2(n-1)}\displaystyle\int_{\Omega}X(R)|\nabla u|^2(u-\varphi)dv\nonumber\\
&&-\dfrac{1}{n(n-1)}\displaystyle\int_{\Omega}u(u-\varphi)X(R)dv
-2\displaystyle\int_{\Omega}u(u-\varphi)\langle\nabla^2u,\mathring{\nabla}^{2}V\rangle dv.
\end{eqnarray}
\end{proposition}

\begin{proof} 
The proof is a long calculation.
Initially, observe that
\[
\operatorname{div}\bigl(Vu\nabla P-P\nabla(Vu)\bigr)
=
Vu\,\Delta P-P\,\Delta(Vu).
\]
Integrating this identity over \(\Omega\) and applying the divergence theorem, we obtain
\begin{equation}\label{eqin}
-\int_{\Omega}Vu\,\Delta Pdv
=
-\int_{\partial\Omega}Vu\,P_{\nu}dv
+\int_{\partial\Omega}P\,(Vu)_{\nu}dv
-\int_{\Omega}P\,\Delta(Vu)dv.
\end{equation}

Next, using the equation \(\Delta u=-(1+nku)\) together with Lemma \ref{lemma1}, we compute
\begin{align}\label{eqnablaVu}
\Delta(Vu)
&=
V\Delta u+u\Delta V+2\langle\nabla V,\nabla u\rangle \nonumber\\
&=
-V(1+nku)-\frac{uX(R)}{2(n-1)}-\frac{R}{n-1}Vu
+2\langle\nabla V,\nabla u\rangle \nonumber\\
&=
-V-\frac{Vu}{n-1}\bigl(R+n(n-1)k\bigr)
+2\langle\nabla V,\nabla u\rangle
-\frac{uX(R)}{2(n-1)}.
\end{align}

Substituting \eqref{eqnablaVu} into \eqref{eqin}, we arrive at
\begin{align}\label{eq43}
-\int_{\Omega}Vu\,\Delta Pdv
={}&
-\int_{\partial\Omega}Vu\,P_{\nu}d\sigma
+\int_{\partial\Omega}P\,(Vu)_{\nu}d\sigma
+\int_{\Omega}PVdv \nonumber\\
&\quad
+\frac{1}{n-1}\int_{\Omega}PVu\bigl(R+n(n-1)k\bigr)dv
-2\int_{\Omega}P\langle \nabla V,\nabla u\rangle dv\nonumber\\
&\quad
+\frac{1}{2(n-1)}\int_{\Omega}Pu\,X(R)dv.
\end{align}

We now make use of the auxiliary function \(\varphi\). Applying integration by parts, we obtain
\begin{align}\label{eq44}
\int_{\partial\Omega}P(Vu)_{\nu}d\sigma
&=
\int_{\partial\Omega}P\bigl(V(u-\varphi)\bigr)_{\nu}d\sigma
+
\int_{\partial\Omega}P(\varphi V)_{\nu}d\sigma
\nonumber\\
&=
\int_{\partial\Omega}PV(u_{\nu}-\varphi_{\nu})d\sigma
+
\int_{\partial\Omega}Pu\,V_{\nu}d\sigma
+
\int_{\partial\Omega}PV\varphi_{\nu}d\sigma
\nonumber\\
&=
\int_{\partial\Omega}PV(u_{\nu}-\varphi_{\nu})d\sigma
+
\int_{\Omega}\operatorname{div}(uP\nabla V)dv
+
\int_{\partial\Omega}PV\varphi_{\nu}d\sigma
\nonumber\\
&=
\int_{\partial\Omega}PV(u_{\nu}-\varphi_{\nu})d\sigma
+
\int_{\Omega}P\langle\nabla u,\nabla V\rangle dv
+
\int_{\Omega}u\langle\nabla P,\nabla V\rangle dv
\nonumber\\
&\quad
+
\int_{\Omega}uP\,\Delta Vdv
+
\int_{\partial\Omega}PV\varphi_{\nu}d\sigma
\nonumber\\
&=
\int_{\partial\Omega}PV(u_{\nu}-\varphi_{\nu})d\sigma
+
\int_{\Omega}P\langle\nabla u,\nabla V\rangle dv
+
\int_{\Omega}u\langle\nabla P,\nabla V\rangle dv
\nonumber\\
&\quad
-
\frac{1}{n-1}
\int_{\Omega}uP\left(\frac{X(R)}{2}+VR\right)dv
+
\int_{\partial\Omega}PV\varphi_{\nu}d\sigma.
\end{align}

Next, integrating by parts once again, we obtain
\begin{align}\label{eq45}
\int_{\partial\Omega}VP\varphi_{\nu}d\sigma
&=
\int_{\Omega}\operatorname{div}(VP\nabla\varphi)dv
\nonumber\\
&=
\int_{\Omega}VP\,\Delta\varphi dv
+
\int_{\Omega}P\langle\nabla V,\nabla\varphi\rangle dv
+
\int_{\Omega}V\langle\nabla P,\nabla\varphi\rangle dv
\nonumber\\
&=
-\int_{\Omega}VPdv
-nk\int_{\Omega}VP\varphi dv
+
\int_{\Omega}P\langle\nabla V,\nabla\varphi\rangle dv
+
\int_{\Omega}V\langle\nabla P,\nabla\varphi\rangle dv,
\end{align}
where we used the identity \(\Delta\varphi=-(1+nk\varphi)\). On the other hand, recalling that the \(P\)-function is given by
\begin{equation}\label{P_func1}
P=|\nabla u|^{2}+\frac{2}{n}u+ku^{2},    
\end{equation}
we compute
\begin{equation}\label{P_funct2}
\nabla P
=
2\nabla^{2}u(\nabla u,\cdot)
+\frac{2}{n}\nabla u
+2ku\nabla u.    
\end{equation}
Substituting these expressions into \eqref{eq45}, we arrive at
\begin{align}\label{eq46}
\int_{\partial\Omega}VP\varphi_{\nu} d\sigma
&=
-\int_{\Omega}VP dv
-nk\int_{\Omega}VP\varphi dv
+
\int_{\Omega}P\langle\nabla V,\nabla\varphi\rangle dv
\nonumber\\
&\quad
+
2\int_{\Omega}V\nabla^{2}u(\nabla u,\nabla\varphi)dv
+
\frac{2}{n}\int_{\Omega}V\langle\nabla u,\nabla\varphi\rangle dv
\nonumber\\
&\quad
+
2k\int_{\Omega}uV\langle\nabla u,\nabla\varphi\rangle dv.
\end{align}

In order to handle the Hessian term appearing in \eqref{eq46}, we recall the Ricci identity
\begin{equation}\label{Ricci_id}
\operatorname{div}(\nabla^2u)
=
\nabla(\Delta u)+\operatorname{Ric}(\nabla u,\cdot).    
\end{equation}
Using this formula, we compute
\begin{align}\label{eq_auxiliar}
\operatorname{div}\bigl(Vu\nabla_{\nabla\varphi}\nabla u\bigr)
&=
\operatorname{div}\bigl(Vu\,\nabla^{2}u(\nabla\varphi,\cdot)\bigr)
\nonumber\\
&=
u\nabla^{2}u(\nabla V,\nabla\varphi)
+
V\langle\nabla u,\nabla_{\nabla\varphi}\nabla u\rangle
+
Vu\,\operatorname{div}(\nabla_{\nabla\varphi}\nabla u)
\nonumber\\
&=
u\nabla^{2}u(\nabla V,\nabla\varphi)
+
V\nabla^{2}u(\nabla u,\nabla\varphi)
+
Vu\,\operatorname{Ric}(\nabla u,\nabla\varphi)
\nonumber\\
&\quad
+
Vu\langle\nabla\Delta u,\nabla\varphi\rangle
+
Vu\langle\nabla^{2}u,\nabla^{2}\varphi\rangle.
\end{align}

Integrating \eqref{eq_auxiliar} over \(\Omega\) and using the divergence theorem, we deduce
\begin{align}\label{eq47}
\int_{\Omega}V\nabla^{2}u(\nabla u,\nabla\varphi)dv
&=
\int_{\partial\Omega}Vu\,\nabla^{2}u(\nabla\varphi,\nu)d\sigma
-
\int_{\Omega}u\nabla^{2}u(\nabla V,\nabla\varphi)dv
\nonumber\\
&\quad
-
\int_{\Omega}Vu(\operatorname{Ric}-nkg)(\nabla u,\nabla\varphi)dv
-
\int_{\Omega}Vu\langle\nabla^{2}u,\nabla^{2}\varphi\rangle dv,
\end{align}
where we have used \(\nabla\Delta u=-nk\nabla u\).

Since \(\varphi\) and \(u\) satisfy \eqref{deltaphi1} and \eqref{serrinoka}, respectively, we infer that
\begin{align}
\int_{\Omega}Vu\langle\nabla^{2}u,\nabla^{2}\varphi\rangle dv
&=
\int_{\Omega}Vu\left[-\left(\frac{1}{n}+k\varphi\right)\Delta u\right]dv
\nonumber\\
&=
\frac{1}{n}\int_{\Omega}Vu\,\Delta\varphi\,(-1-nku)dv
\nonumber\\
&=
-\frac{1}{n}\int_{\Omega}Vu\,\Delta\varphi dv
-k\int_{\Omega}Vu^{2}\Delta\varphi dv.
\nonumber
\end{align}

Applying the divergence theorem once again, together with the identities
\[
\operatorname{div}(Vu\nabla\varphi)
=
u\langle\nabla V,\nabla\varphi\rangle
+
V\langle\nabla u,\nabla\varphi\rangle
+
Vu\Delta\varphi,
\]
and
\[
\operatorname{div}(Vu^{2}\nabla\varphi)
=
\langle\nabla(Vu^{2}),\nabla\varphi\rangle
+
Vu^{2}\Delta\varphi,
\]
we obtain
\begin{align}
\int_{\Omega}Vu\langle\nabla^{2}u,\nabla^{2}\varphi\rangle dv
&=
\frac{1}{n}\int_{\Omega}u\langle\nabla V,\nabla\varphi\rangle dv
+
\frac{1}{n}\int_{\Omega}V\langle\nabla u,\nabla\varphi\rangle dv
\nonumber\\
&\quad
-
\frac{1}{n}\int_{\partial\Omega}Vu\,\varphi_{\nu}d\sigma
+
k\int_{\Omega}\langle\nabla(Vu^{2}),\nabla\varphi\rangle dv
-
k\int_{\partial\Omega}Vu^{2}\varphi_{\nu}d\sigma.
\nonumber
\end{align}

Proceeding further, we arrive at
\begin{align}\label{eq48}
\int_{\Omega}Vu\langle\nabla^{2}u,\nabla^{2}\varphi\rangle dv
&=
-\frac{1}{n}\int_{\partial\Omega}Vu\,\varphi_{\nu}d\sigma
+
\frac{1}{n}\int_{\Omega}u\langle\nabla V,\nabla\varphi\rangle dv
+
\frac{1}{n}\int_{\Omega}V\langle\nabla u,\nabla\varphi\rangle dv
\nonumber\\
&\quad
+
2k\int_{\Omega}Vu\langle\nabla u,\nabla\varphi\rangle dv
+
k\int_{\Omega}u^{2}\langle\nabla V,\nabla\varphi\rangle dv
-
k\int_{\partial\Omega}Vu^{2}\varphi_{\nu}d\sigma.
\end{align}

Plugging this into \eqref{eq47}, we obtain
\begin{align}\label{eqstar}
\int_{\Omega}V\nabla^{2}u(\nabla u,\nabla\varphi)dv
&=
\int_{\partial\Omega}Vu\,\nabla^{2}u(\nabla\varphi,\nu)d\sigma
-
\int_{\Omega}u\nabla^{2}u(\nabla V,\nabla\varphi)dv
\nonumber\\
&\quad
-
\int_{\Omega}Vu(\operatorname{Ric}-nkg)(\nabla u,\nabla\varphi)dv
+
\frac{1}{n}\int_{\partial\Omega}Vu\,\varphi_{\nu}d\sigma
\nonumber\\
&\quad
-
\frac{1}{n}\int_{\Omega}u\langle\nabla V,\nabla\varphi\rangle dv
-
\frac{1}{n}\int_{\Omega}V\langle\nabla u,\nabla\varphi\rangle dv
\nonumber\\
&\quad
+
k\int_{\partial\Omega}Vu^{2}\varphi_{\nu}d\sigma
-
2k\int_{\Omega}Vu\langle\nabla u,\nabla\varphi\rangle dv
\nonumber\\
&\quad
-
k\int_{\Omega}u^{2}\langle\nabla V,\nabla\varphi\rangle dv.
\end{align}

Substituting \eqref{eqstar} into \eqref{eq46}, we obtain
\begin{align}\label{eqtwostar}
\int_{\partial\Omega}VP\varphi_{\nu}d\sigma
&=
-\int_{\Omega}VPdv
-nk\int_{\Omega}VP\varphi dv
+
\int_{\Omega}P\langle\nabla V,\nabla\varphi\rangle dv
+
\frac{2}{n}\int_{\Omega}V\langle\nabla u,\nabla\varphi\rangle dv
\nonumber\\
&\quad
+
2k\int_{\Omega}uV\langle\nabla u,\nabla\varphi\rangle dv
+
2\int_{\partial\Omega}Vu\nabla^{2}u(\nabla\varphi,\nu) d\sigma
\nonumber\\
&\quad
-
2\int_{\Omega}u\nabla^{2}u(\nabla V,\nabla\varphi)dv
-
2\int_{\Omega}Vu(\operatorname{Ric}-nkg)(\nabla u,\nabla\varphi)dv
\nonumber\\
&\quad
+
\frac{2}{n}\int_{\partial\Omega}Vu\,\varphi_{\nu}d\sigma
-
\frac{2}{n}\int_{\Omega}u\langle\nabla V,\nabla\varphi\rangle dv
-
\frac{2}{n}\int_{\Omega}V\langle\nabla u,\nabla\varphi\rangle dv
\nonumber\\
&\quad
+
2k\int_{\partial\Omega}Vu^{2}\varphi_{\nu}d\sigma
-
4k\int_{\Omega}Vu\langle\nabla u,\nabla\varphi\rangle dv
\nonumber\\
&\quad
-
2k\int_{\Omega}u^{2}\langle\nabla V,\nabla\varphi\rangle dv.
\end{align}

After collecting similar terms, this becomes
\begin{align}
\int_{\partial\Omega}VP\varphi_{\nu} d\sigma
&=
-\int_{\Omega}VP dv
+
\int_{\Omega}P\langle\nabla V,\nabla\varphi\rangle dv
+
2\int_{\partial\Omega}Vu\nabla^{2}u(\nabla\varphi,\nu) d\sigma
\nonumber\\
&\quad
-
2k\int_{\Omega}u^{2}\langle\nabla V,\nabla\varphi\rangle dv
-
2\int_{\Omega}Vu(\operatorname{Ric}-(n-1)kg)(\nabla u,\nabla\varphi) dv
\nonumber\\
&\quad
-
2\int_{\Omega}u\nabla^{2}u(\nabla V,\nabla\varphi) dv
-
\frac{2}{n}\int_{\Omega}u\langle\nabla V,\nabla\varphi\rangle dv
\nonumber\\
&\quad
+
\frac{2}{n}\int_{\partial\Omega}Vu\,\varphi_{\nu} d\sigma
+
2k\int_{\partial\Omega}Vu^{2}\varphi_{\nu} d\sigma
-
nk\int_{\Omega}VP\varphi dv.
\label{eqtwostar_simplified}
\end{align}

Substituting \eqref{eqtwostar_simplified} into \eqref{eq44}, we obtain
\begin{align}\label{eqstar44}
\int_{\partial\Omega}P(Vu)_{\nu}d\sigma
&=
\int_{\partial\Omega}PV(u_{\nu}-\varphi_{\nu})d\sigma
+
\int_{\Omega}P\langle\nabla u,\nabla V\rangle dv
+
\int_{\Omega}u\langle\nabla P,\nabla V\rangle dv
\nonumber\\
&\quad
-
\frac{1}{n-1}
\int_{\Omega}uP\left(\frac{1}{2}X(R)+VR\right)dv
-
\int_{\Omega}VP dv
\nonumber\\
&\quad
+
\int_{\Omega}P\langle\nabla V,\nabla\varphi\rangle dv
+
2\int_{\partial\Omega}Vu\nabla^{2}u(\nabla\varphi,\nu) d\sigma
\nonumber\\
&\quad
-
2k\int_{\Omega}u^{2}\langle\nabla V,\nabla\varphi\rangle dv
-
2\int_{\Omega}Vu(\operatorname{Ric}-(n-1)kg)(\nabla u,\nabla\varphi) dv
\nonumber\\
&\quad
-
2\int_{\Omega}u\nabla^{2}u(\nabla V,\nabla\varphi)dv
-
\frac{2}{n}\int_{\Omega}u\langle\nabla V,\nabla\varphi\rangle dv
\nonumber\\
&\quad
+
\frac{2}{n}\int_{\partial\Omega}Vu\,\varphi_{\nu}d\sigma
+
2k\int_{\partial\Omega}Vu^{2}\varphi_{\nu}d\sigma
-
nk\int_{\Omega}VP\varphi dv.
\end{align}

Returning now to \eqref{eq43}, we obtain
\begin{align}\label{eq49}
-\int_{\Omega}Vu\,\Delta P dv
&=
-\int_{\partial\Omega}Vu\,P_{\nu}d\sigma
+
\int_{\partial\Omega}VP(u_{\nu}-\varphi_{\nu})d\sigma
-
\int_{\Omega}P\langle\nabla V,\nabla(u-\varphi)\rangle dv
\nonumber\\
&\quad
+
\int_{\Omega}u\langle\nabla V,\nabla P\rangle dv
+
2\int_{\partial\Omega}Vu\nabla^{2}u(\nabla\varphi,\nu) d\sigma
\nonumber\\
&\quad
-
2\int_{\Omega}u\nabla^{2}u(\nabla V,\nabla\varphi) dv
-
\frac{2}{n}\int_{\Omega}u\langle\nabla V,\nabla\varphi\rangle dv
\nonumber\\
&\quad
+
\frac{2}{n}\int_{\partial\Omega}Vu\,\varphi_{\nu} d\sigma
+
2k\int_{\partial\Omega}Vu^{2}\varphi_{\nu} d\sigma
\nonumber\\
&\quad
-
2k\int_{\Omega}u^{2}\langle\nabla V,\nabla\varphi\rangle dv
+
nk\int_{\Omega}VP(u-\varphi) dv
\nonumber\\
&\quad
-
2\int_{\Omega}Vu(\operatorname{Ric}-(n-1)kg)(\nabla u,\nabla\varphi) dv.
\end{align}

Next, we collect the terms in \eqref{eq49} involving integrals over \(\Omega\) that do not contain \(k\). Denoting their sum by
\[
C
=
-\int_{\Omega}P\langle\nabla V,\nabla(u-\varphi)\rangle dv
+
\int_{\Omega}u\langle\nabla V,\nabla P\rangle dv
-
2\int_{\Omega}u\nabla^{2}u(\nabla V,\nabla\varphi)dv
-
\frac{2}{n}\int_{\Omega}u\langle\nabla V,\nabla\varphi\rangle dv,
\]
and expanding the expression of \(P\), see \eqref{P_func1} and \eqref{P_funct2}, we find
\begin{align}\label{eq410}
C
&=
-\int_{\Omega}|\nabla u|^{2}
\langle\nabla V,\nabla(u-\varphi)\rangle dv
+
2\int_{\Omega}u\nabla^{2}u\bigl(\nabla V,\nabla(u-\varphi)\bigr) dv
\nonumber\\
&\quad
+
k\int_{\Omega}u^{2}\langle\nabla V,\nabla u\rangle dv
+
k\int_{\Omega}u^{2}\langle\nabla V,\nabla\varphi\rangle dv.
\end{align}

We now compute the first term appearing in \eqref{eq410}. To begin with, observe that
\[
\operatorname{div}\bigl(-|\nabla u|^{2}(u-\varphi)\nabla V\bigr)
=
-(u-\varphi)\langle\nabla |\nabla u|^{2},\nabla V\rangle
-
|\nabla u|^{2}\langle\nabla(u-\varphi),\nabla V\rangle
-
|\nabla u|^{2}(u-\varphi)\Delta V.
\]

Integrating over \(\Omega\) and applying the divergence theorem, we obtain
\begin{align}\label{eqaux0}
-\int_{\Omega}|\nabla u|^{2}\langle\nabla V,\nabla(u-\varphi)\rangle dv
&=
-\int_{\partial\Omega}|\nabla u|^{2}(u-\varphi)V_{\nu} d\sigma
+
\int_{\Omega}(u-\varphi)\langle\nabla |\nabla u|^{2},\nabla V\rangle dv
\nonumber\\
&\quad
+
\int_{\Omega}|\nabla u|^{2}(u-\varphi)\Delta V dv.
\end{align}
Using the identity \eqref{Lapl_V}, we conclude that
\begin{align}\label{eqaux0_final}
-\int_{\Omega}|\nabla u|^{2}\langle\nabla V,\nabla(u-\varphi)\rangle dv
&=
-\int_{\partial\Omega}|\nabla u|^{2}(u-\varphi)V_{\nu} d\sigma
+
2\int_{\Omega}(u-\varphi)\nabla^2 u(\nabla u,\nabla V) dv
\nonumber\\
&\quad
-
\frac{1}{2(n-1)}
\int_{\Omega}X(R)|\nabla u|^{2}(u-\varphi)dv
\nonumber\\
&\quad
-
\frac{1}{n-1}
\int_{\Omega}|\nabla u|^{2}(u-\varphi)VR dv.
\end{align}

On the other hand, using once again the Ricci identity \eqref{Ricci_id}, we obtain
\begin{align*}
\operatorname{div}\bigl(u(u-\varphi)\nabla^{2}u(\nabla V)\bigr)
&=
(u-\varphi)\nabla^{2}u(\nabla V,\nabla u)
+
u\nabla^{2}u(\nabla V,\nabla(u-\varphi))
\nonumber\\
&\quad
+
u(u-\varphi)(\operatorname{Ric}-nkg)(\nabla V,\nabla u)
\nonumber\\
&\quad
+
u(u-\varphi)\langle\nabla^{2}u,\nabla^{2}V\rangle.
\end{align*}

Integrating over \(\Omega\) and applying the divergence theorem, we deduce
\begin{align}\label{auxeq1}
\int_{\Omega}u\nabla^{2}u(\nabla V,\nabla(u-\varphi)) dv
&=
\int_{\partial\Omega}u(u-\varphi)\nabla^{2}u(\nabla V,\nu)d\sigma
-
\int_{\Omega}(u-\varphi)\nabla^{2}u(\nabla V,\nabla u)dv
\nonumber\\
&\quad
-
\int_{\Omega}u(u-\varphi)(\operatorname{Ric}-nkg)(\nabla V,\nabla u)dv
-
\int_{\Omega}u(u-\varphi)\langle\nabla^{2}u,\nabla^{2}V\rangle dv.
\end{align}

Since
\[
\nabla^{2}V
=
\mathring{\nabla^{2}V}
+
\frac{\Delta V}{n}\,g,
\]
using \eqref{Lapl_V} we infer that
\begin{align}\label{eqaux1}
\int_{\Omega}u\nabla^{2}u(\nabla V,\nabla(u-\varphi))dv
&=
\int_{\partial\Omega}u(u-\varphi)\nabla^{2}u(\nabla V,\nu)d\sigma
-
\int_{\Omega}(u-\varphi)\nabla^{2}u(\nabla V,\nabla u) dv
\nonumber\\
&\quad
-
\int_{\Omega}u(u-\varphi)(\operatorname{Ric}-nkg)(\nabla V,\nabla u) dv
-
\int_{\Omega}u(u-\varphi)
\langle\nabla^{2}u,\mathring{\nabla^{2}V}\rangle dv
\nonumber\\
&\quad
-
\frac{1}{2n(n-1)}
\int_{\Omega}u(u-\varphi)X(R) dv
-
\frac{1}{n(n-1)}
\int_{\Omega}u(u-\varphi)VR dv
\nonumber\\
&\quad
-
\frac{k}{2(n-1)}
\int_{\Omega}u^{2}(u-\varphi)X(R) dv
-
\frac{k}{n-1}
\int_{\Omega}u^{2}(u-\varphi)VR dv.
\end{align}

Combining \eqref{eq410}, \eqref{eqaux0_final} and \eqref{eqaux1}, we obtain
\begin{align*}
C
&=
-\int_{\partial\Omega}|\nabla u|^{2}(u-\varphi)V_{\nu}d\sigma
+2\int_{\partial\Omega}u(u-\varphi)\nabla^{2}u(\nabla V,\nu)d\sigma
\nonumber\\
&\quad
-2\int_{\Omega}u(u-\varphi)(\operatorname{Ric}-nkg)(\nabla u,\nabla V)dv
-2\int_{\Omega}u(u-\varphi)\langle\nabla^{2}u,\mathring{\nabla^{2}V}\rangle dv
\nonumber\\
&\quad
-\frac{1}{n-1}\int_{\Omega}(u-\varphi)
\left(
\frac{2}{n}uVR
+2kVu^{2}R
+|\nabla u|^{2}VR
\right)dv
\nonumber\\
&\quad
-\frac{1}{n-1}\int_{\Omega}(u-\varphi)
\left(
\frac12 X(R)|\nabla u|^{2}
+\frac{1}{n}uX(R)
+ku^{2}X(R)
\right)dv
+k\int_{\Omega}u^{2}\langle\nabla V,\nabla(u+\varphi)\rangle dv.
\end{align*}

Substituting this expression for \(C\) into \eqref{eq49}, we deduce
\begin{align}\label{eq411}
-\int_{\Omega}Vu\,\Delta P dv
&=
-\int_{\partial\Omega}|\nabla u|^{2}(u-\varphi)V_{\nu} d\sigma
+2\int_{\partial\Omega}u(u-\varphi)\nabla^{2}u(\nabla V,\nu)d\sigma
\nonumber\\
&\quad
-\int_{\partial\Omega}Vu\,P_{\nu}d\sigma
+\int_{\partial\Omega}VP(u_{\nu}-\varphi_{\nu})d\sigma
+2\int_{\partial\Omega}Vu\nabla^{2}u(\nabla\varphi,\nu)d\sigma
\nonumber\\
&\quad
-\frac{2}{n(n-1)}
\int_{\Omega}u(u-\varphi)VR dv
-\frac{2k}{n-1}
\int_{\Omega}Vu^{2}(u-\varphi)R dv
\nonumber\\
&\quad
-2\int_{\Omega}u(u-\varphi)(\operatorname{Ric}-nkg)(\nabla u,\nabla V)dv
+k\int_{\Omega}u^{2}\langle\nabla V,\nabla(u-\varphi)\rangle dv
\nonumber\\
&\quad
-\frac{1}{n-1}
\int_{\Omega}|\nabla u|^{2}(u-\varphi)VR dv
-\frac{1}{2(n-1)}
\int_{\Omega}X(R)|\nabla u|^{2}(u-\varphi)dv
\nonumber\\
&\quad
-\frac{1}{n(n-1)}
\int_{\Omega}u(u-\varphi)X(R)dv
-\frac{k}{n-1}
\int_{\Omega}u^{2}(u-\varphi)X(R) dv
\nonumber\\
&\quad
-2\int_{\Omega}u(u-\varphi)
\langle\nabla^{2}u,\mathring{\nabla^{2}V}\rangle dv
+\frac{2}{n}\int_{\partial\Omega}Vu\,\varphi_{\nu}d\sigma
+2k\int_{\partial\Omega}Vu^{2}\varphi_{\nu}d\sigma
\nonumber\\
&\quad
+nk\int_{\Omega}VP(u-\varphi)dv
-2\int_{\Omega}Vu(\operatorname{Ric}-(n-1)kg)(\nabla u,\nabla\varphi)dv.
\end{align}

We now expand \(P\) and \(P_{\nu}\) in the previous identity to obtain
\begin{align}\label{eq412}
-\int_{\Omega}Vu\,\Delta P dv
&=
-\int_{\partial\Omega}|\nabla u|^{2}(u-\varphi)V_{\nu} d\sigma
+2\int_{\partial\Omega}u(u-\varphi)\nabla^{2}u(\nabla V,\nu) d\sigma
\nonumber\\
&\quad
+\int_{\partial\Omega}V|\nabla u|^{2}(u_{\nu}-\varphi_{\nu}) d\sigma
-2\int_{\partial\Omega}Vu\nabla^{2}u(\nabla(u-\varphi),\nu) d\sigma
\nonumber\\
&\quad
-k\int_{\partial\Omega}Vu^{2}(u-\varphi)_{\nu} d\sigma
-\frac{2}{n(n-1)}
\int_{\Omega}Vu(u-\varphi)(R-n(n-1)k)dv
\nonumber\\
&\quad
+k\int_{\Omega}u^{2}\langle\nabla V,\nabla(u-\varphi)\rangle dv
-2\int_{\Omega}u(u-\varphi)(\operatorname{Ric}-(n-1)kg)(\nabla u,\nabla V)dv
\nonumber\\
&\quad
-2\int_{\Omega}Vu(\operatorname{Ric}-(n-1)kg)(\nabla u,\nabla\varphi) dv
+2k\int_{\Omega}u(u-\varphi)\langle\nabla u,\nabla V\rangle dv
\nonumber\\
&\quad
-\frac{1}{n-1}
\int_{\Omega}V|\nabla u|^{2}(u-\varphi)(R-(n-1)nk) dv
-\frac{2k}{n-1}
\int_{\Omega}Vu^{2}(u-\varphi)
\left(R-\frac{(n-1)nk}{2}\right) dv
\nonumber\\
&\quad
-\frac{1}{2(n-1)}
\int_{\Omega}X(R)|\nabla u|^{2}(u-\varphi)dv
-\frac{1}{n(n-1)}
\int_{\Omega}u(u-\varphi)X(R)dv
\nonumber\\
&\quad
-\frac{k}{n-1}
\int_{\Omega}u^{2}(u-\varphi)X(R)dv
-2\int_{\Omega}u(u-\varphi)
\langle\nabla^{2}u,\mathring{\nabla^{2}V}\rangle dv.
\end{align}

Next, observe that
\[
\operatorname{div}\bigl((u-\varphi)u^{2}\nabla V\bigr)
=
u^{2}\langle\nabla(u-\varphi),\nabla V\rangle
+
(u-\varphi)\langle\nabla u^{2},\nabla V\rangle
+
(u-\varphi)u^{2}\Delta V.
\]
Integrating by parts, we obtain
\begin{align}\label{eqaux4}
k\int_{\Omega}u(u-\varphi)\langle\nabla V,\nabla u\rangle dv
&=
\frac{k}{2}
\int_{\Omega}(u-\varphi)\langle\nabla V,\nabla u^{2}\rangle dv
\nonumber\\
&=
\frac{k}{2}
\Bigg[
\int_{\partial\Omega}(u-\varphi)u^{2}V_{\nu}d\sigma
-
\int_{\Omega}u^{2}\langle\nabla(u-\varphi),\nabla V\rangle dv
-\int_{\Omega}(u-\varphi)u^{2}\Delta V dv
\Bigg]
\nonumber\\
&=
\frac{k}{2}
\Bigg[
\int_{\partial\Omega}(u-\varphi)u^{2}V_{\nu} d\sigma
+
\frac{1}{2(n-1)}
\int_{\Omega}(u-\varphi)u^{2}X(R) dv
\nonumber\\
&\qquad\qquad
+
\frac{1}{n-1}
\int_{\Omega}(u-\varphi)u^{2}VR dv
-
\int_{\Omega}u^{2}\langle\nabla(u-\varphi),\nabla V\rangle dv
\Bigg],
\end{align}
where we used \eqref{Lapl_V}. Therefore,
\begin{align}\label{eq413}
-\int_{\Omega}Vu\,\Delta P dv
&=
\int_{\partial\Omega}|\nabla u|^{2}
\bigl[V(u-\varphi)_{\nu}-(u-\varphi)V_{\nu}\bigr] d\sigma
+
2\int_{\partial\Omega}
u\nabla^{2}u\bigl((u-\varphi)\nabla V-V\nabla(u-\varphi),\nu\bigr)d\sigma
\nonumber\\
&\quad
+
k\int_{\partial\Omega}
u^{2}\bigl[(u-\varphi)V_{\nu}-V(u-\varphi)_{\nu}\bigr]d\sigma
-\frac{2}{n(n-1)}
\int_{\Omega}Vu(u-\varphi)(R-n(n-1)k) dv
\nonumber\\
&\quad
-2\int_{\Omega}u(u-\varphi)
(\operatorname{Ric}-(n-1)kg)(\nabla u,\nabla V)dv
-2\int_{\Omega}Vu
(\operatorname{Ric}-(n-1)kg)(\nabla u,\nabla\varphi)dv
\nonumber\\
&\quad
-\frac{k}{n-1}
\int_{\Omega}Vu^{2}(u-\varphi)(R-(n-1)nk) dv
-\frac{1}{n-1}
\int_{\Omega}V|\nabla u|^{2}(u-\varphi)(R-(n-1)nk)dv
\nonumber\\
&\quad
-\frac{k}{2(n-1)}
\int_{\Omega}u^{2}(u-\varphi)X(R)dv
-\frac{1}{2(n-1)}
\int_{\Omega}X(R)|\nabla u|^{2}(u-\varphi)dv
\nonumber\\
&\quad
-\frac{1}{n(n-1)}
\int_{\Omega}u(u-\varphi)X(R)dv
-2\int_{\Omega}u(u-\varphi)
\langle\nabla^{2}u,\mathring{\nabla^{2}V}\rangle dv.
\end{align}

This completes the proof of \eqref{eqlema}.
\end{proof}

\begin{remark}\label{remark_main}
The identity in Proposition~\ref{intlemma} is particularly relevant to the study of the stability of spherical symmetry in Riemannian manifolds, in contexts related both to Alexandrov's soap bubble theorem and to Serrin's overdetermined problems, through the use of integral identity techniques; see \cite{poggesi1, poggesi, OO21, SZ25}. It also plays an important role in the analysis of partially overdetermined problems in space forms, where Serrin-type conditions are used to characterize the rigidity of umbilical hypersurfaces (see \cite{guoxia1, guoxia2}). Both applications rely on fundamental integral identities that are of independent interest; see, for instance, \cite[Theorem 2.1]{poggesi} and \cite[Lemma 4.2]{xie2025}.    
\end{remark}

In particular, we highlight the form taken by the last result in the special setting of Einstein manifolds.

\begin{proof}[Proof of Theorem \ref{MainTh_2}]
Since $\operatorname{Ric}=(n-1)kg$, all terms involving
$R-n(n-1)k$ and $\operatorname{Ric}-(n-1)kg$ in Proposition~\ref{intlemma}
vanish identically. Moreover, since Einstein manifolds have constant scalar curvature, we also have $X(R)=0$. Finally, every boundary term containing $u$ vanishes because $u=0$ on $\partial\Omega$.

It remains to analyze the boundary contribution
\[
\int_{\partial\Omega}|\nabla u|^{2}
\bigl[V(u-\varphi)_{\nu}-(u-\varphi)V_{\nu}\bigr]d\sigma.
\]
Since $u=0$ on $\partial\Omega$, we have
$|\nabla u|^{2}=u_{\nu}^{2}$ along $\partial\Omega$, and therefore
\[
\int_{\partial\Omega}|\nabla u|^{2}
\bigl[V(u-\varphi)_{\nu}-(u-\varphi)V_{\nu}\bigr]d\sigma
=
\int_{\partial\Omega}u_{\nu}^{2}
\bigl[V(u-\varphi)_{\nu}-(u-\varphi)V_{\nu}\bigr]d\sigma.
\]

Now, for any constant $c^{2}$, integration by parts yields
\begin{align*}
c^{2}\int_{\partial\Omega}
\bigl[V(u-\varphi)_{\nu}-(u-\varphi)V_{\nu}\bigr]d\sigma
&=
c^{2}\int_{\Omega}
\Bigl(V\Delta(u-\varphi)-(u-\varphi)\Delta V\Bigr)d\sigma.
\end{align*}

Using $\Delta u=-1-nku$, $\Delta\varphi=-1-nk\varphi$, and $\Delta V=-nkV$, which follows from \eqref{Lapl_V}, we obtain
\begin{align*}
c^{2}\int_{\partial\Omega}
\bigl[V(u-\varphi)_{\nu}-(u-\varphi)V_{\nu}\bigr]d\sigma
&=
c^{2}\int_{\Omega}
\Bigl(V(-nk(u-\varphi))+nkV(u-\varphi)\Bigr)d\sigma \\
&=0.
\end{align*}

Consequently,
\begin{align*}
\int_{\partial\Omega}u_{\nu}^{2}
\bigl[V(u-\varphi)_{\nu}-(u-\varphi)V_{\nu}\bigr]d\sigma
&=
\int_{\partial\Omega}(u_{\nu}^{2}-c^{2})
\bigl[V(u-\varphi)_{\nu}-(u-\varphi)V_{\nu}\bigr]d\sigma.
\end{align*}
Finally, by Lemma~\ref{lemma_tracefree}, we conclude that $\mathring{\nabla}^{2}V=0$, and hence the last integral term in Proposition~\ref{intlemma} vanishes. Therefore, since $\Delta P = 2|\mathring{\nabla}^{2}u|^{2}$, the desired identity follows.



% Finally, the inequalities in {\rm(i)} and {\rm(ii)} follow immediately from the sign of the term
% \[
% -2\int_{\Omega}u(u-\varphi)
% \langle \nabla^{2}u,\mathring{\nabla}^{2}V\rangle.
% \]
\end{proof}
\begin{remark}\label{existence}
In the setting of Theorem \ref{MainTh_2}, it follows from \eqref{Lapl_V} that
\[
\Delta V=-nkV.
\]
Therefore, if $V$ does not change sign, the problem \eqref{Dirichlet2} can be rewritten as
\begin{equation}\label{Dirichlet*}
\Delta u -\frac{\Delta V}{V} u = -1 \quad \text{in } \Omega, 
\qquad 
u = 0 \quad \text{on } \partial \Omega.
\end{equation}
The existence and uniqueness of solutions to \eqref{Dirichlet*} are guaranteed, for example, by \cite[Theorem 2.3]{foga}; see also \cite[Lemma 2.5]{li-xia}.
\end{remark}
In the next result, we show that, in conjunction with the overdetermined boundary condition \(u_{\nu}=-c\), the previous theorem immediately yields a rigidity statement. Although this result was proved via the \(P\)-function method in \cite[Corollary 1]{andrade2025} (see also \cite[Theorem 2]{freitas2024} for the closed conformal case), our approach follows directly from the integral identity established above.

\begin{corollary}\label{cor:rigidity}
Let \( (M^n, g) \) be an Einstein manifold with $\operatorname{Ric} = (n-1)k g,$ for some \( k \in \mathbb{R} \). Assume additionally that \( M \) is endowed with a conformal vector field \( X \) whose conformal factor does not change sign. If \( u \) is a solution of the overdetermined problem
\[
\begin{cases}
\Delta u + nku = -1 & \text{in } \Omega,\\
u>0 & \text{in } \Omega,\\
u=0 & \text{on } \partial\Omega,\\
u_{\nu}=-c & \text{on } \partial\Omega,
\end{cases}
\]
for some bounded domain \( \Omega\subset M \) with \(C^{2}\)-boundary and some constant \(c>0\), then \( \Omega \) is a metric ball and \( u \) is radial.
\end{corollary}

% Let $(M^n, g)$ be an Einstein Riemannian manifold, i.e. $\operatorname{Ric}=(n-1)Kg$, $K\in\mathbb{R}$, endowed with a conformal vector field $X$ with conformal factor $V$ that does not change sign ($V>0$ and $u\leq \varphi$ or $V<0$ and $u \geq \varphi$). If $\Omega$ is a bounded domain with $C^2$-boundary and $u$ is a solution of \eqref{serrinoka} and \eqref{serrinokb} with $u>0$ in $\Omega$, then $\Omega$ is a metric ball and $u$ is a radial function.

\begin{proof}
By Theorem~\ref{MainTh_2}, we have
\[
-2\int_{\Omega}Vu\,|\mathring{\nabla}^{2}u|^{2}dv=0.
\]
Since \(V\) does not change sign and \(u>0\) in \(\Omega\), it follows that $\mathring{\nabla}^{2}u=0$ in $\Omega$. The conclusion then follows directly from \cite[Lemma 6]{farina2022}.
\end{proof}

% \begin{proof}
% Suppose that $V>0$ and $u\leq \varphi$ in $\Omega$. Then
% \begin{eqnarray}\label{VUDeltaP}
% -Vu\Delta P\leq 0\quad \text{in } \Omega.
% \end{eqnarray}
% Since $|\nabla u|^2=c^2$ on $\partial \Omega$ and $\operatorname{Ric}=(n-1)kg$, by Lemma \ref{intlemma} we obtain
% \begin{eqnarray*}
% 0\geq \displaystyle\int_{\Omega}-Vu\Delta P
% &=&c^2\displaystyle\int_{\partial \Omega}[V(u-\varphi)_{\nu}-(u-\varphi)V_{\nu}]\\
% &=&\displaystyle\int_{\Omega}(V\Delta(u-\varphi)-(u-\varphi)\Delta V)
% -2\displaystyle\int_{\Omega}u(u-\varphi)\langle\nabla^2u, \mathring{\nabla^2V}\rangle \geq 0,
% \end{eqnarray*}
% where we have used that $\Delta u=-(1+nKu)$, $\Delta \varphi=-(1+nK\varphi)$ and $u\leq \varphi$. This implies that
% \[
% \int_{\Omega}-Vu\Delta P=0.
% \]
% Together with \eqref{VUDeltaP}, this implies that $\Delta P =0$ in $\Omega$ since $V>0$ and $u>0$ in $\Omega$. Finally, using Lemma \ref{farina-roncoroni} we finish the proof.
% \end{proof}










\section*{Data availability statement}
This manuscript has no associated data.

\section*{Conflict of interest statement}
On behalf of all authors, the corresponding author states that there is no conflict of interest.

\section*{Acknowledgments}
M. Andrade was partially supported by the Brazilian National Council for Scientific and Technological Development (CNPq, grants 408834/2023-4, 403869/2024-2, and 400078/2025-2) and FAPITEC/SE/Brazil (grant 019203.01303/2024-1). She also extends her gratitude to the Department of Mathematics at Princeton University, where this work was completed while she was a Visiting Fellow, for its hospitality. She is especially thankful to Ana Menezes for her warm hospitality and continuous encouragement.

A. Freitas was partially supported by the Brazilian National Council for Scientific and Technological Development (CNPq), grants 308141/2025-3 and 406078/2025-4. He also wishes to thank IMECC/UNICAMP for the stimulating scientific environment and support provided during the development of part of this work. In particular, he is grateful to Lino Grama for his warm hospitality and constant encouragement.

The authors would like to thank Márcio Santos for his discussions on the subject of this paper and valuable suggestions.
%\bibliography{main.bib}
%\bibliographystyle{acm}
\begin{thebibliography}{99}

\bibitem{andrade2025}
M. Andrade, A. Freitas, and D. A. Mar\'in,
\emph{Rigidity results for Serrin's overdetermined problems in Riemannian manifolds.}
Calc. Var. Partial Differential Equations {\bf 64} (2025), no. 4, 119.

\bibitem{ciraolo}
G. Ciraolo and L. Vezzoni,
\emph{On Serrin’s overdetermined problem in space forms.}
Manuscripta Math. {\bf 159} (2019), 445--452.

\bibitem{farina2022}
A. Farina and A. Roncoroni,
\emph{Serrin's type problems in warped product manifolds.}
Commun. Contemp. Math. {\bf 24} (2022), no. 4, 2150020.

\bibitem{foga}
M. Fogagnolo and A. Pinamonti,
\emph{New integral estimates in substatic Riemannian manifolds and the Alexandrov theorem.}
J. Math. Pures Appl. {\bf 163} (2022), 299--317.

\bibitem{freitas2024}
A. Freitas, A. Roncoroni, and M. Santos,
\emph{A note on Serrin's type problem on Riemannian manifolds.}
J. Geom. Anal. {\bf 34} (2024), no. 7, 200.

\bibitem{guoxia1}
J. Guo and C. Xia,
\emph{A partially overdetermined problem in a half ball.}
Calc. Var. Partial Differential Equations {\bf 58} (2019), no. 5, 160.

\bibitem{guoxia2}
J. Guo and C. Xia,
\emph{A partially overdetermined problem in domains with partial umbilical boundary in space forms.}
Adv. Calc. Var. {\bf 17} (2024), no. 1, 11--31.

\bibitem{kr}
W. Kuhnel and H. B. Rademacher,
\emph{Conformal vector fields on pseudo-Riemannian spaces.}
Differential Geom. Appl. {\bf 7} (1997), 237--255.

\bibitem{li-xia}
J. Li and C. Xia,
\emph{An integral formula and its applications on sub-static manifolds.}
J. Differential Geom. {\bf 113} (2019), 493--518.

\bibitem{poggesi1}
R. Magnanini and G. Poggesi,
\emph{On the stability for Alexandrov’s Soap Bubble Theorem.}
J. Anal. Math. {\bf 139} (2019), 179--205.

\bibitem{poggesi}
R. Magnanini and G. Poggesi,
\emph{Serrin's Problem and Alexandrov's Soap Bubble Theorem: Enhanced Stability via Integral Identities.}
Indiana Univ. Math. J. {\bf 69} (2020), no. 4, 1181--1205.

\bibitem{markvorsen2006}
S. Markvorsen and V. Palmer,
\emph{Torsional rigidity of minimal submanifolds.}
Proc. London Math. Soc. {\bf 93} (2006), 253--272.

\bibitem{montiel}
S. Montiel,
\emph{Unicity of constant mean curvature hypersurfaces in some Riemannian manifolds.}
Indiana Univ. Math. J. {\bf 48} (1999), no. 2, 711--748.

\bibitem{obata}
M. Obata,
\emph{Certain conditions for a Riemannian manifold to be isometric with a sphere.}
J. Math. Soc. Japan {\bf 14} (1962), 333--340.

\bibitem{OO21}
Y. Okamoto and M. Onodera,
\emph{Stability analysis of an overdetermined fourth order boundary value problem via an integral identity.}
J. Differential Equations {\bf 301} (2021), 97--111.

\bibitem{petersen}
P. Petersen,
\emph{Riemannian Geometry.}
Springer, New York, 2006.

\bibitem{pohozaev1965}
S. I. Poho\v{z}aev,
\emph{On the eigenfunctions of the equation $\Delta u + \lambda f(u)=0$.}
Dokl. Akad. Nauk SSSR {\bf 165} (1965), 36--39.

\bibitem{polya1951}
G. P\'olya and G. Szeg\"o,
\emph{Isoperimetric Inequalities in Mathematical Physics.}
Ann. of Math. Stud., vol. 27, Princeton Univ. Press, Princeton, NJ, 1951.

\bibitem{reilly}
R. Reilly,
\emph{Applications of the Hessian operator in a Riemannian manifold.}
Indiana Univ. Math. J. {\bf 26} (1977), no. 3, 459--472.

\bibitem{Roncoroni2018}
A. Roncoroni,
\emph{A Serrin-type symmetry result on model manifolds: an extension of the Weinberger argument.}
C. R. Math. Acad. Sci. Paris {\bf 356} (2018), no. 6, 648--656.

\bibitem{SZ25}
J. Scheuer and X. Zhang,
\emph{Stability of the Wulff shape with respect to anisotropic curvature functionals.}
J. Funct. Anal. {\bf 288} (2025), no. 3, 110715.

\bibitem{serrin1971}
J. Serrin,
\emph{A symmetry problem in potential theory.}
Arch. Rational Mech. Anal. {\bf 43} (1971), 304--318.

\bibitem{sokolniko1956}
I. S. Sokolnikoff,
\emph{Mathematical Theory of Elasticity.}
McGraw--Hill, New York, 1956.

\bibitem{struwe}
M. Struwe,
\emph{Variational Methods: Applications to Nonlinear Partial Differential Equations and Hamiltonian Systems.}
Ergeb. Math. Grenzgeb. (3), Springer-Verlag, Berlin--Heidelberg, 1990.

\bibitem{vandenberg1990}
M. van den Berg,
\emph{Large time asymptotics of the heat flow.}
Quart. J. Math. Oxford {\bf 41} (1990), 245--253.

\bibitem{wang2020}
Q. Wang and C. Xia,
\emph{The Pohozaev identity for the anisotropic $p$-Laplacian and estimates of the torsion function.}
Rev. Mat. Iberoam. {\bf 36} (2020), no. 4, 1257--1278.

\bibitem{xie2025}
Y. Xie,
\emph{Characterizations of umbilical hypersurfaces by partially overdetermined problems in space forms.}
J. Math. Study {\bf 58} (2025), no. 3, 286--313.

\bibitem{yano}
K. Yano,
\emph{The Theory of Lie Derivatives and its Applications.}
North-Holland, Amsterdam, 1957.

\end{thebibliography}
\end{document}
